\documentclass[12pt,a4paper,oneside]{book}

% PAQUETES DE CONFIGURACIÓN BÁSICA
\usepackage[spanish]{babel}
\selectlanguage{spanish}
\usepackage[utf8]{inputenc}
\usepackage[T1]{fontenc} 
\usepackage{mathptmx} % Fuente tipo Times
\renewcommand{\baselinestretch}{1.3} % Interlineado
\usepackage{hyperref}
\usepackage{bookmark}
% Mejora el manejo de comprobaciones de archivos auxiliares y reruns necesarios
\usepackage{rerunfilecheck}

% GEOMETRÍA Y MÁRGENES
\usepackage{geometry}
\geometry{left=4cm, right=2.5cm, top=3.5cm, bottom=2.5cm}

% PAQUETES PARA CITAS (Estilo APA)
\usepackage{apacite}
\let\cite\shortcite 

% PAQUETES PARA GRÁFICOS, TABLAS Y OTROS
\usepackage{graphicx}
\usepackage{amsmath, amssymb, amsfonts}
\usepackage{url}
\usepackage{ragged2e} % Para \begin {justify}
\usepackage{booktabs}
\usepackage{longtable}
\usepackage{pdfpages} % para insertar páginas PDF externas

% ENCABEZADOS Y PIES DE PÁGINA
\usepackage{fancyhdr}
\pagestyle{fancy}
\fancyhf{} % Limpia todo
\fancyhead[R]{\thepage} % Número página arriba-derecha
\renewcommand{\headrulewidth}{0pt} % Sin línea encabezado
\setlength{\headheight}{14.5pt} % Soluciona advertencia de fancyhdr
\addtolength{\topmargin}{-2.5pt} % Compensa el aumento de headheight

% RENOMBRAR TÍTULOS
\renewcommand{\chaptername}{Capítulo}
\renewcommand{\contentsname}{Índice}
\renewcommand{\listtablename}{Índice de Cuadros}
\renewcommand{\listfigurename}{Índice de Figuras}

% COMIENZA EL DOCUMENTO
\begin{document}

% PORTADA
\begin{titlepage}
    \centering
    \vspace*{1cm}
    
\includegraphics[width=4cm]{logofacen.jpg}\par
    \vspace{1cm}
    {\large\bfseries UNIVERSIDAD NACIONAL DE ASUNCIÓN}\par
    {\large Facultad de Ciencias Exactas y Naturales}\par
    Departamento de Estadística\par
    Maestría en Estadística\par
    \vspace{2cm}
    {\Large\bfseries Análisis de datos de panel del balance hídrico en Paraguay: Cuantificación de la sequía y la aridificación ante el calentamiento global (2000–2023)}\par % TÍTULO DE TESIS
    \vspace{2cm}

    {\large\bfseries Hernán J. Vargas Ojeda}\par
    \vskip 0.5 cm

    \begin{tabular}{l}
    \hskip 5 cm    Orientador: Prof. Dr. Gonzalo Martin Vicente\\
    \hskip 5 cm   Co-Orientador: Prof. Dr. Gustavo Ignacio Rivas M.\\
    \end{tabular}
    
    \vfill % Empuja el texto hacia abajo
    Tesis presentada a la Facultad de Ciencias Exactas y Naturales,
    Universidad Nacional de Asunción, como requisito para la 
    obtención del Grado de Magíster en Estadística\par
    \vspace{2cm}
    SAN LORENZO - PARAGUAY\par
    Junio - 2026\par % AÑO DE DEFENSA
\end{titlepage}

% PÁGINA DE AUTORIDADES / FICHA 
\newpage
\thispagestyle{empty} % Esta página no tendrá número
\pagenumbering{roman} % A partir de aquí, la numeración es en números romanos
\setcounter{page}{2} % Empezamos a contar desde ii

\textbf{Ficha Catalográfica del Trabajo de Tesis}
\vspace{1cm}

\begin{justify} 
Datos Internacionales de Catalogación en la Publicación (CIP) \\
DE LA BIBLIOTECA E INTERNET DE LA FACEN - UNA
\end{justify}
\vspace{2cm}

\fbox{\parbox[b]{0.9\linewidth}{ 
    \vspace{0.1cm}
    \textit{Vargas Ojeda, Hernán Javier}\vspace{0.2cm}\\
     Análisis de datos de panel del balance hídrico en Paraguay: Cuantificación de la sequía y la aridificación ante el calentamiento global (2000–2023) / Hernán J. Vargas Ojeda. – San Lorenzo: Universidad Nacional de Asunción. Facultad de Ciencias Exactas y Naturales, 2026.\vspace{0.2cm}\\ 
    % número de páginas final;
    78 h.; 30 cm.\vspace{0.2cm}\\
    Incluye referencias bibliográficas.\vspace{0.2cm}\\
    Tesis (Magíster en Estadística). -- Universidad Nacional de Asunción. Facultad de Ciencias Exactas y Naturales, 2026.\vspace{0.2cm}\\ 
    % Palabras clave relevantes de la tesis
    1. Tesis de postgrado – UNA 2. Sequía 3. SPEI 4. Aridificación\\
    5. Datos de panel 6. Calentamiento Global 7. Paraguay\\
    \vspace{0.2cm}
    \begin{flushright}
        % Clasificación de la biblioteca
        519.5/V297a 
    \end{flushright}
}}
\clearpage % Salto de página

% PÁGINA DE APROBACIÓN / MESA EXAMINADORA
% pendiente

% DEDICATORIA
\chapter*{Dedicatoria}
\addcontentsline{toc}{chapter}{Dedicatoria} % Añade la dedicatoria al índice
\vspace*{10cm} % espacio vertical 
\begin{flushright}
    \textit{A mi familia...} %dedicatoria aquí
\end{flushright}
\clearpage

% AGRADECIMIENTOS
\chapter*{Agradecimientos}
\addcontentsline{toc}{chapter}{Agradecimientos} % Añade al índice
En primer lugar, deseo expresar mi más profundo agradecimiento a mis orientadores, por su invaluable guía y paciencia infinita. Sus agudas observaciones y su constante apoyo fueron fundamentales para dar forma y dirección a esta investigación.

Agradezco a los miembros de la mesa examinadora de Coordinación de Postgrado e Investigación en Estadística,
Departamento de Estadística, FACEN - UNA, por aceptar evaluar este trabajo y por sus valiosas sugerencias que sin duda lo enriquecerán.

Mi gratitud a la Facultad de Ciencias Exactas y Naturales de la Universidad Nacional de Asunción, y en especial al Departamento de Estadística, por brindarme el espacio y los recursos para mi formación de posgrado.

Quisiera extender un agradecimiento especial a la Dirección de Meteorología e Hidrología (DMH) por facilitar el acceso a los datos crudos que constituyen la base de este estudio.

A mis compañeros de la Maestría en Estadística, gracias por las largas jornadas de estudio, las discusiones enriquecedoras y el compañerismo que hicieron este camino más ameno.

A mis padres, María Asunción de Rodríguez y Mario Rodríguez M., por su amor incondicional y por inculcarme el valor del esfuerzo y la educación. A mi pareja, Adriana Viveros González, por su paciencia infinita, su aliento en los momentos difíciles y por ser mi principal apoyo. Sin ustedes, este logro no habría sido posible.
\clearpage

% RESUMEN
\chapter*{Resumen}
\addcontentsline{toc}{chapter}{Resumen} % Añade al índice
\vspace{1cm}
La presente investigación evalúa la evolución temporal y espacial de las sequías y los procesos de aridificación en Paraguay durante el período 2000--2023, analizando el impacto del factor térmico y la demanda evaporativa en un contexto de calentamiento global. Con este propósito, se estructuró una base de datos en formato de datos de panel balanceado a partir de los registros diarios de la red oficial de bases meteorológicas oficiales del país. El territorio nacional fue regionalizado mediante técnicas de agrupamiento jerárquico aglomerativo (HAC) y optimizado mediante el coeficiente de silueta promedio, identificando dos subregiones climáticas homogéneas de balance hídrico. Las estimaciones de tendencias estacionales y regionales se realizaron mediante un modelo de regresión de datos de panel de efectos aleatorios, utilizando el test de Hausman para su validación estadística y aplicando los errores estándar de Driscoll y Kraay (1998) robustos a la correlación serial y a la fuerte dependencia espacial confirmada. Los resultados revelan una marcada divergencia estacional en el balance hídrico: mientras las estaciones de otoño e invierno se mantuvieron hídricamente estables, las épocas cálidas de primavera y verano registraron trayectorias de aridificación. La tasa de aridificación estival más intensa del país se concentró en la subregión agro productiva del sureste. El marco de contraste paralelo entre el modelo de balance climático (SPEI-6) y el de precipitación univariada (SPI-6) demostró la validez del rol del estrés térmico, evidencia que las sequías en Paraguay están adquiriendo una naturaleza predominantemente térmica en el siglo XXI, donde la evaporación forzada por el incremento sistemático de las temperaturas máximas extremas induce procesos de aridificación estacional incluso bajo volúmenes de precipitación estables.
\vspace{1cm}

\noindent\textbf{Palabras clave:} Datos de panel, balance hídrico, aridificación, SPEI, efectos aleatorios, Driscoll-Kraay, Calentamiento global, Paraguay.
\clearpage

% ABSTRACT (SUMMARY)
\chapter*{Abstract} % Título en inglés
\addcontentsline{toc}{chapter}{Abstract} % Añade al índice
\vspace{1cm}
This research evaluates the temporal and spatial evolution of droughts and seasonal aridification processes in Paraguay during the 2000--2023 period, analyzing the impact of the thermal factor and evaporative demand under global warming. To achieve this, a balanced panel database was structured utilizing daily records from the country's official meteorological station network. The national territory was regionalized using hierarchical agglomerative clustering (HAC) and optimized via the average silhouette width, which identified two water-balance homogeneous subregions. Seasonal and regional trends were estimated using a random effects panel data regression model, validated statistically through the Hausman specification test, and adjusted using Driscoll and Kraay (1998) robust standard errors to control for serial autocorrelation and the strong cross-sectional dependence confirmed. The results reveal a marked seasonal divergence in the water balance: while autumn and winter seasons remained hydrologically stable, spring and summer warm seasons registered drying trends. The most intense summer aridification rate in the country was concentrated in the agricultural heartland of the South-East subregion, evidence that droughts in Paraguay are acquiring a predominantly thermal nature in the 21st century, where evaporation forced by the systematic increase in extreme maximum temperatures induces seasonal aridification processes even under stable precipitation volumes.
\vspace{1cm}

\noindent\textbf{Keywords:} Panel data, water balance, aridification, SPEI, random effects, Driscoll-Kraay, Global warming, Paraguay.
\clearpage

% ÍNDICES
\renewcommand{\contentsname}{Índice}
\tableofcontents
\clearpage

\listoftables
\addcontentsline{toc}{chapter}{\listtablename} % Añade al índice principal
\clearpage

\listoffigures
\addcontentsline{toc}{chapter}{\listfigurename} % Añade al índice principal
\clearpage

% CUERPO PRINCIPAL DE LA TESIS
\pagenumbering{arabic} % Empieza la numeración arábiga

\chapter{INTRODUCCIÓN}

En todo balance hay un ingreso y un egreso, y en el \textit{balance hídrico} los ingresos son las lluvias y los egresos están representados por la evaporación potencial (ETP), que básicamente es la cantidad de agua que el aire intenta extraer del suelo y las plantas. Cuando este balance tiene un déficit estamos ante una \textit{sequía}, situación que se corrige cuando vuelven las lluvias, por lo que tradicionalmente se pensaba que las sequías eran causadas únicamente por la falta de lluvias. 

La intensificación de las sequías altera estructuralmente el balance hídrico local, fuerza al ecosistema hacia un nuevo equilibrio más seco convirtiendo una anomalía climática en un proceso de \textit{aridificación} estructural, el balance hídrico se inclina hacia un déficit sistemático que redefine el clima local a largo plazo.

En décadas recientes, se ha descubierto que altas temperaturas aumenta la ETP convirtiendose en un actor decisivo que puede forzar un déficit hídrico incluso cuando llueve lo normal \cite{vicenteserrano2010}. La temperatura desempeña un rol fundamental en la dinámica de eventos climáticos, según \citeA{massondelmotte2021}, el incremento de las temperaturas globales provocado por el calentamiento global acelera la tasa de evaporación del agua del suelo y la transpiración de las plantas.

El Grupo Intergubernamental de Expertos sobre el Cambio Climático (IPCC) sostiene de manera concluyente que el calentamiento del sistema climático global es inequívoco \cite{stocker_climate_2013}. Este \textit{calentamiento global}, impulsado por el aumento de gases de efecto invernadero de origen antropogénico, no solo altera los patrones de precipitación, sino que también eleva las temperaturas, incrementando la evapotranspiración y acelerando la pérdida de humedad del suelo, de esta manera, el calentamiento global actúa como un multiplicador del riesgo de sequía, haciéndola más probable, intensa y persistente \cite{massondelmotte2021}.

\section{Planteamiento del problema}

La literatura nacional reciente ha documentado tendencias ascendentes y estadísticamente significativas en las temperaturas medias anuales en gran parte del territorio nacional \cite{chavez2023}, así como alteraciones en los regímenes de precipitación y en la frecuencia de eventos climáticos extremos \cite{chavez2023, grassi2020}. Sin embargo, aún no se han abordado de manera sistemática la sequía y los procesos de aridificación en el Paraguay bajo el contexto del calentamiento global. El problema fundamental radica en que este vacío de conocimiento genera una incertidumbre climática que se traduce en riesgos no cuantificados para la seguridad hídrica y la estabilidad económica del país.

Esta investigación aborda dicha problemática bajo un enfoque multivariado, incorporando la temperatura en el análisis del balance hídrico. El modelo propuesto considera el rol térmico bajo dos dimensiones esenciales: por un lado, su vínculo directo con el aumento de la demanda evaporativa de la atmósfera, lo cual incide sobre la frecuencia, intensidad, duración y extensión espacial de los eventos de sequía; por otro lado, analiza su comportamiento estacional intrínseco. Históricamente, el régimen climático de Paraguay se ha caracterizado por veranos húmedos e inviernos relativamente más secos; no obstante, esta configuración estacional clásica podría estar experimentando modificaciones estructurales como consecuencia del calentamiento global.

A partir de esta brecha de conocimiento en Paraguay, se formulan las siguientes preguntas de investigación considerando el período de estudio 2000--2023:
\begin{enumerate}
    \item ¿Cómo se manifiestan las sequías en Paraguay en términos de su frecuencia, intensidad, duración y extensión espacial? 
    \item ¿Existen evidencias estadísticamente significativas de procesos de aridificación en Paraguay, y cómo varía su intensidad entre estaciones climáticas y regiones geográficas? 
    \item ¿En qué medida la integración de la temperatura y la precipitación en el análisis del balance hídrico permite detectar condiciones significativamente más deficitarias, atribuibles al aumento de las temperaturas asociadas al calentamiento global?
\end{enumerate}

\section{Justificación de la investigación}

El desarrollo de este trabajo se justifica bajo tres dimensiones analíticas:
\begin{itemize}
    \item \textit{Dimensión socioeconómica:} contar con una métrica precisa de un balance hídrico y sus  dinámicas provee evidencia empírica clave para el diseño de políticas de adaptación sectorial, optimizando la gestión de riesgos en sectores altamente vulnerables como el agropecuario.
    \item \textit{Dimensión científica:} supera las limitaciones de análisis climáticos previos al proveer una cuantificación que permita una caracterización sistemática de la sequía en Paraguay incorporando un enfoque multivariado que integra la temperatura como un factor condicionante, mejorando su comprensión en el marco del calentamiento global.
    \item \textit{Dimensión metodológica:} introduce el cálculo de un índice integrado de balance hídrico sensible al calentamiento global y un modelado de datos de panel no explorado para el contexto nacional, proporcionando un marco de análisis replicable y adaptable para futuras investigaciones sobre sequía y aridificación en la región.
\end{itemize}

\section{Hipótesis}

Para guiar el análisis empírico, se plantean las siguientes hipótesis de investigación:
\begin{itemize}
    \item \textit{H1}: Las sequías en Paraguay durante el período 2000--2023 muestran diferencias significativas en su frecuencia, severidad, duración y extensión de los eventos.
    \item \textit{H2}: El período 2000--2023 evidencia procesos de aridificación estadísticamente significativos en Paraguay con una intensidad variable entre estaciones climáticas y regiones geográficas.
    \item \textit{H3}: La integración de la temperatura y la precipitación en la evaluación del balance hídrico permite detectar condiciones significativamente más deficitarias que aquellas estimadas únicamente por la precipitación, reflejando el efecto del incremento de las temperaturas asociado al calentamiento global.
\end{itemize}

\section{Objetivos de la investigación}

\subsection{Objetivo general}

Evaluar la evolución temporal del balance hídrico en Paraguay (2000--2023), mediante el análisis integrado de la precipitación y la temperatura que nos permita, carácterizar los eventos de sequía y ajustar un modelo de datos de panel para detectar dinámicas de largo plazo en un contexto de calentamiento global.

\subsection{Objetivos específicos}
\begin{enumerate}
    \item Construir una base de datos climática en formato de panel para Paraguay durante el período 2000--2023, integrando registros mensuales de precipitación y temperatura sometidos a controles de consistencia climática.
    \item Estimar las series mensuales del índice de precipitación - evapotranspiración estandarizado (SPEI) como variable de estudio del déficit hídrico.
    \item Regionalizar el territorio paraguayo a partir de la variabilidad común del balance hídrico, mediante técnicas de agrupamiento espacial.
    \item Cuantificar las tendencias temporales del SPEI mediante modelos de datos de panel, a fin de identificar y evaluar los procesos de aridificación según su variación espacial y estacional.
    \item Evaluar la contribución del incremento de las temperaturas en el déficit hídrico, mediante un análisis comparativo y estadístico entre las dinámicas del SPEI y el índice estandarizado de precipitación (SPI).
\end{enumerate}
\chapter{Marco Teórico}

\section{La Sequía: Características}

La sequía es un fenómeno de desarrollo lento y acumulativo. Sus efectos se manifiestan de manera gradual, lo que dificulta la determinación precisa de su inicio y final, pudiendo persistir durante años incluso después de que las condiciones meteorológicas hayan retornado a la normalidad \cite{mishra2010}.

Es fundamental distinguir el concepto de sequía del de aridez. La aridez es una característica estructural, climatológica y permanente de una región, definida por niveles de precipitación sistemáticamente bajos. En contraste, la sequía representa una anomalía temporal; es decir, un déficit hídrico transitorio con respecto a las condiciones normales esperadas para esa región. Por lo tanto, un evento de sequía puede manifestarse tanto en regiones húmedas como en regiones áridas. 

De acuerdo con la literatura especializada, las cuatro características fundamentales para el análisis de la sequía son la intensidad, la duración, la frecuencia y la extensión espacial \cite{mishra2010}. El análisis conjunto de estas cuatro dimensiones resulta esencial para una evaluación integral del riesgo, así como para el diseño de sistemas de alerta temprana y la formulación de estrategias de mitigación efectivas.

\section{Clasificación de la Sequía}

La comunidad científica ha desarrollado un marco conceptual estándar que clasifica la sequía en cuatro tipos principales, entendidos como una secuencia de propagación de impactos \cite{mishra2010}:
\begin{enumerate}
    \item \textbf{Sequía meteorológica:} Constituye el punto de partida del fenómeno y se define estrictamente por un déficit de precipitación durante un período prolongado en comparación con la media histórica registrada en una región.
    \item \textbf{Sequía agrícola:} Se manifiesta cuando el déficit de lluvia se propaga induciendo un déficit de humedad en la zona radicular del suelo, afectando la germinación, el crecimiento y el rendimiento de los cultivos, así como la productividad de las pasturas destinadas a la ganadería.
    \item \textbf{Sequía hidrológica:} Refiere a la disminución significativa de los niveles y caudales en las fuentes de agua superficiales y subterráneas, tales como ríos, lagos, embalses y acuíferos.
    \item \textbf{Sequía socioeconómica:} Ocurre cuando la escasez física del recurso hídrico, derivada de las fases anteriores, afecta de manera directa la oferta y la demanda de bienes y servicios económicos.
\end{enumerate}

La presente investigación se centra en la caracterización y modelización de la \textit{sequía meteorológica}. Esta elección se fundamenta en que representa el indicador primario y precursor de todos los demás impactos hidrológicos, agrícolas y socioeconómicos, convirtiéndose en el componente crítico para cualquier sistema de detección temprana. Asimismo, desde una perspectiva instrumental, la sequía meteorológica se define de manera directa a partir de las variables físicas disponibles en la red de estaciones terrestres (precipitación y temperatura), evitando los requerimientos de datos biofísicos complejos y de difícil acceso que demandan las otras clasificaciones.

\section{Causas y Consecuencias de la Sequía en Paraguay}

Al igual que otros países de la región, el régimen pluviométrico de Paraguay se encuentra modulado por fenómenos climáticos de gran escala. El principal motor de la variabilidad interanual es El Niño-Oscilación del Sur (ENSO), un patrón de calentamiento y enfriamiento anómalo de las aguas superficiales del Océano Pacífico Ecuatorial. 

Además del ENSO, otros patrones de variabilidad climática a nivel hemisférico, como la Oscilación Decadal del Pacífico (PDO) y la Oscilación Antártica (AAO), influyen en la regularidad de las precipitaciones, aunque sus dinámicas operan en escalas de tiempo más prolongadas y complejas \cite{vera2006}. Superpuesta a esta variabilidad climática natural, la tendencia de fondo del calentamiento global antropogénico actúa como un factor crítico de intensificación. El incremento sostenido de las temperaturas eleva la evapotranspiración potencial (ETP), exacerbando la severidad de las sequías bajo un mecanismo conocido en la literatura como "hotter droughts" \cite{massondelmotte2021}. Este proceso fue determinante en la reciente sequía extrema registrada en la Cuenca del Plata, donde las altas temperaturas extremaron el déficit hídrico preexistente provocado por la fase de La Niña \cite{naumann2023}.

Dada la estructura productiva de su economía y su posición geográfica, Paraguay presenta una alta susceptibilidad ante las consecuencias de la sequía. Estos impactos se manifiestan de manera transversal en los siguientes sectores clave:
\begin{itemize}
    \item \textbf{Sector agropecuario:} La sequía agrícola reduce drásticamente el rendimiento de los cultivos de renta de mayor relevancia para el producto interno bruto (como la soja y el maíz), disminuye la receptividad de las pasturas y puede inducir tasas de mortandad ganadera significativas, traduciéndose en pérdidas económicas millonarias \cite{cepal2014}.
    \item \textbf{Navegación fluvial:} El estiaje prolongado de los ríos Paraguay y Paraná, consecuencia directa de la sequía hidrológica, afecta la principal vía de importación y exportación comercial de la nación. La reducción del calado operativo de las barcazas incrementa sustancialmente los costos logísticos y debilita la competitividad del comercio exterior.
    \item \textbf{Generación de energía hidroeléctrica:} Paraguay depende casi en su totalidad de la hidroelectricidad generada por las represas binacionales de Itaipú y Yacyretá. La sequía hidrológica reduce el caudal afluente, comprometiendo la capacidad de generación y disminuyendo los ingresos estatales percibidos por la exportación de excedentes energéticos.
    \item \textbf{Aspectos sociales y ambientales:} La escasez de agua afecta de manera directa a las comunidades rurales y suburbanas vulnerables, limitando su acceso al agua potable y comprometiendo la agricultura de subsistencia. Adicionalmente, las condiciones prolongadas de sequedad y calor elevan exponencialmente el riesgo de propagación de incendios forestales, con la consecuente degradación de ecosistemas y pérdida de biodiversidad.
\end{itemize}

\section{Consistencia y Completitud en Series Temporales Climatológicas}

Un desafío recurrente en la modelización climatológica es la presencia de registros faltantes. El tratamiento inapropiado de la falta de datos mediante la eliminación sistemática de periodos observacionales introduce un sesgo que disminuye sustancialmente el poder estadístico de los modelos y altera el cálculo de acumulaciones temporales consecutivas esenciales para los índices de sequía.

De acuerdo con las directrices metodológicas de la Organización Meteorológica Mundial \cite{wmo2017}, la imputación de huecos cortos en series temporales es una práctica estandarizada y necesaria para asegurar la homogeneidad y consistencia de los datos. El tratamiento de registros faltantes en esta investigación se fundamenta teóricamente en la preservación de la estructura de autocorrelación de las series continuas. Para variables con una alta inercia térmica temporal, como la temperatura, se emplean interpolaciones lineales sobre ventanas temporales restringidas. En el caso de variables intermitentes y asimétricas, como la precipitación, la imputación se restringe a evaluaciones contextuales directas de vecindad temporal. Al limitar de manera estricta el proceso de imputación a huecos de duración ultra-corta (1 a 2 días), se evita la inserción de sesgo o variabilidad artificial, garantizando la consistencia analítica requerida para el posterior cálculo multiescalar de los índices de sequía.

\section{Cuantificación de la Sequía Meteorológica: Los Índices SPI y SPEI}

Dada la naturaleza compleja y de desarrollo acumulativo de la sequía, su monitoreo y modelización rigurosa requieren el uso de indicadores cuantitativos estandarizados. Estos índices sintetizan las fluctuaciones climáticas instrumentales en un único valor numérico, permitiendo comparar la severidad espacial y temporal de las anomalías independientemente de la climatología local.

\subsection{El Índice de Precipitación Estandarizada (SPI)}

El Índice de Precipitación Estandarizada (SPI), desarrollado por McKee, Doesken y Kleist \cite{mckee1993}, es uno de los indicadores más implementados globalmente. El SPI parte de la premisa matemática de que el déficit acumulativo de precipitación representa el principal factor detrás de las condiciones de sequía. 

Su procedimiento metodológico consiste en ajustar los registros históricos de precipitación acumulada para una escala temporal seleccionada (por ejemplo, 3, 6 o 12 meses) a una función de distribución de probabilidad teórica, típicamente la distribución Gamma. Posteriormente, las probabilidades acumuladas resultantes se transforman mediante la inversa de la función de distribución normal estándar, produciendo un índice normalizado ($Z$-score) con media cero y varianza unitaria. Los valores negativos indican condiciones más secas de lo normal y los valores positivos condiciones húmedas. Su gran ventaja radica en su flexibilidad multiescalar y simplicidad operacional, al requerir únicamente registros de precipitación.

A pesar de su utilidad regulatoria, el SPI exhibe una limitación conceptual crítica en escenarios de cambio climático: ignora sistemáticamente la influencia de las anomalías de temperatura. El calentamiento global altera no solo los regímenes de precipitación, sino que incrementa de forma directa la demanda evaporativa de la atmósfera, un factor físico determinante en el balance hídrico del suelo. La evapotranspiración potencial (ETP) define la pérdida máxima de agua que una superficie vegetal experimentaría si la disponibilidad hídrica fuera ilimitada. Temperaturas más elevadas conducen a una ETP incremental, acelerando la pérdida de humedad de los suelos y la vegetación. Por consiguiente, omitir la dinámica de la ETP puede resultar en una subestimación sustancial de la severidad de la sequía, especialmente durante eventos de olas de calor persistentes \cite{massondelmotte2021}.

\subsection{El Índice de Precipitación-Evapotranspiración Estandarizado (SPEI)}

Para resolver las limitaciones conceptuales del SPI ante el incremento térmico global, Vicente-Serrano, Beguería y López-Moreno \cite{vicenteserrano2010} diseñaron el Índice de Precipitación-Evapotranspiración Estandarizado (SPEI). El SPEI preserva las propiedades matemáticas y la flexibilidad del SPI, pero fundamenta su cálculo en el balance hídrico climático mensual ($D_i$), definido como la diferencia simple entre la precipitación ($P_i$) y la evapotranspiración potencial estimada ($ETP_i$):
\begin{equation}
    D_i = P_i - ETP_i
\end{equation}

El cálculo del índice acumula estas diferencias mensuales en diferentes escalas de tiempo históricas. Debido a que el balance hídrico puede tomar valores negativos, se recurre a la distribución Log-Logística de tres parámetros para ajustar las series temporales, ofreciendo una flexibilidad superior a la distribución Gamma de dos parámetros del SPI. Finalmente, se realiza el proceso de estandarización normal clásico para obtener los valores del SPEI ($Z$-score), permitiendo una clasificación objetiva de la intensidad de las anomalías hídricas (Cuadro 2.1).

\begin{table}[h!]
\centering
\caption{Clasificación de las condiciones de sequía y humedad según los valores del Índice SPEI.}
\label{tab:clasificacion_spei}
\begin{tabular}{ll}
\hline
\textbf{Valor del SPEI} & \textbf{Categoría} \\ \hline
2.00 o más & Extremadamente Húmedo \\
1.50 a 1.99 & Severamente Húmedo \\
1.00 a 1.49 & Moderadamente Húmedo \\
-0.99 a 0.99 & Condiciones Normales \\
-1.00 a -1.49 & Sequía Moderada \\
-1.50 a -1.99 & Sequía Severa \\
-2.00 o menos & Sequía Extrema \\ \hline
\end{tabular}
\end{table}

\subsection{Antecedentes de Investigación sobre la Sequía en Paraguay}

Si bien el estudio de la sequía mediante índices de balance hídrico multivariados es un estándar internacional, su aplicación instrumental en Paraguay ha sido históricamente limitada. No obstante, las investigaciones nacionales previas proveen las bases científicas para examinar las dinámicas térmicas locales.

Un aporte referencial es el de Chávez \cite{chavez2023}, quien mediante pruebas estadísticas no paramétricas aplicadas a series históricas de temperatura (periodo 1960--2018), confirmó tendencias de calentamiento sistemáticas y estadísticamente significativas en la mayor parte de las estaciones meteorológicas terrestres del país. De forma complementaria, Pasten \cite{pasten2007} identificó un incremento significativo en la frecuencia de días y noches cálidas a nivel nacional.

Asimismo, Grassi \cite{grassi2020} determinó que la tasa de calentamiento observada en el Paraguay durante las últimas décadas supera el promedio mundial. Un hallazgo relevante de este autor es que, a pesar de registrarse incrementos leves en las medias de precipitación anual, la frecuencia y los impactos de las sequías extremas han aumentado significativamente durante el siglo XXI. Esto sugiere que las fluctuaciones térmicas positivas desempeñan un rol dominante en la intensificación del déficit hídrico por encima de la dinámica de las precipitaciones. Por último, las proyecciones de la Comisión Económica para América Latina y el Caribe \cite{cepal2014} confirman que la persistencia de estas dinámicas inducirá choques de oferta desfavorables en la estabilidad económica nacional.

A pesar de la relevancia de estos aportes, la literatura climatológica local carece de investigaciones que examinen la evolución del balance hídrico integrado ($P - ETP$) empleando índices estandarizados de última generación como el SPEI bajo un enfoque de modelización de panel.

\section{Fundamentos del Modelado Estadístico para Datos de Panel}

El análisis de series temporales climatológicas observadas a través de redes terrestres de monitoreo presenta un desafío estadístico fundamental: la heterogeneidad espacial no observada. Cada estación meteorológica se halla condicionada por características físicas locales invariantes en el tiempo (como la altitud, la latitud, la topografía circundante o la geomorfología del suelo) que definen su nivel de base o intercepto climatológico.

Si estas características intrínsecas no se controlan de manera adecuada, sus efectos se combinan con la tendencia temporal estimada, induciendo sesgos severos por variable omitida. El análisis de datos de panel, o datos longitudinales, proporciona el marco de inferencia estadístico adecuado para aislar y resolver este problema \cite{wooldridge2010}. Al observar de manera repetida las mismas estaciones meteorológicas a lo largo del tiempo, este enfoque metodológico permite descomponer la varianza observada en dos componentes analíticos fundamentales: la variación entre unidades (variación \textit{between}) y la variación a lo largo del tiempo para una misma unidad (variación \textit{within}).

\subsection{Modelos de Efectos Fijos (FE) vs. Efectos Aleatorios (RE)}

Existen dos aproximaciones teóricas alternativas para especificar estadísticamente la heterogeneidad no observada de las estaciones:

\subsubsection{Modelo de Efectos Fijos (FE)}
Asume que los interceptos de cada unidad de análisis son fijos y específicos para cada estación $i$, representados por el parámetro $\alpha_i$. El estimador de efectos fijos, también conocido como estimador "within", elimina las características invariantes en el tiempo mediante la transformación por diferencias respecto a la media temporal de cada variable para cada unidad. La ventaja principal de este enfoque es que no requiere imponer el supuesto de ortogonalidad entre los efectos fijos específicos y las variables explicativas del modelo, controlando de manera robusta cualquier factor espacial invariante.

La especificación general del modelo FE se define como:
\begin{equation}
    y_{it} = \alpha_i + \beta x_{it} + \varepsilon_{it}
\end{equation}
Donde la transformación que elimina la heterogeneidad no observada opera como:
\begin{equation}
    (y_{it} - \bar{y}_i) = \beta(x_{it} - \bar{x}_i) + (u_{it} - \bar{u}_i)
\end{equation}

\subsubsection{Modelo de Efectos Aleatorios (RE)}
Este modelo asume que las características específicas invariantes de cada estación meteorológica, $\alpha_i$, no actúan como parámetros fijos a estimar, sino como componentes aleatorios distribuidos de manera uniforme dentro del término de error compuesto del sistema. La hipótesis matemática fundamental del modelo de efectos aleatorios reside en la ortogonalidad o ausencia de correlación entre el componente aleatorio individual y los regresores del modelo. Si este supuesto de ortogonalidad se cumple, el estimador RE es estadísticamente más eficiente que el de efectos fijos, al utilizar de forma óptima tanto la variabilidad \textit{within} como la variabilidad \textit{between}.

La ecuación general se plantea de la siguiente forma:
\begin{equation}
    y_{it} = \beta_0 + \beta x_{it} + (\alpha_i + \varepsilon_{it})
\end{equation}
Bajo la condición estricta de ortogonalidad:
\begin{equation}
    Cov(\alpha_i, x_{it}) = 0
\end{equation}

\subsection{Selección del Modelo: El Test de Especificación de Hausman}

La elección metodológica entre la aproximación de efectos fijos y efectos aleatorios no se determina de forma discrecional, sino que se fundamenta en la consistencia estadística del modelo RE. El test de especificación de Hausman \cite{hausman1978} constituye la prueba estadística estándar para evaluar esta idoneidad.

Esta prueba contrasta sistemáticamente el vector de coeficientes estimados por el modelo FE (consistente bajo cualquier supuesto de correlación de los efectos individuales) frente al vector de coeficientes del modelo RE (consistente y eficiente únicamente bajo el supuesto de ortogonalidad). La hipótesis nula del test establece que no existen diferencias sistemáticas significativas entre ambos estimadores:
\begin{equation}
    H_0: \hat{\beta}_{FE} = \hat{\beta}_{RE}
\end{equation}

Si el p-valor calculado para la estadística de prueba es inferior al nivel de significancia establecido (típicamente $\alpha = 0.05$), se rechaza la hipótesis nula, indicando la presencia de correlación y la obligatoriedad metodológica de implementar la aproximación de efectos fijos para garantizar la consistencia de las estimaciones. Por el contrario, un p-valor elevado sugiere que las estimaciones del modelo de efectos aleatorios son robustas y preferibles debido a su mayor eficiencia.

\subsection{Supuestos y Limitaciones del Modelo de Panel}

La validez metodológica del modelo de efectos aleatorios descansa sobre la verificación de los siguientes supuestos analíticos clave:
\begin{enumerate}
    \item \textbf{Exogeneidad estricta:} Se asume que el componente de error idiosincrático, $\varepsilon_{it}$, no se encuentra correlacionado con las variables explicativas para ningún periodo temporal observacional. La consistencia del efecto aleatorio individual $\alpha_i$ respecto a las variables explicativas se valida mediante el test de Hausman.
    \item \textbf{Homocedasticidad y no autocorrelación del error:} Supone que los errores idiosincráticos presentan una varianza constante e independiente a través de las observaciones temporales para una misma entidad. Debido a que las series climatológicas suelen exhibir autocorrelación y heterocedasticidad, este estudio aborda preventivamente dicha potencial desviación mediante la estimación de errores estándar robustos agrupados (clustered robust standard errors) a nivel de estación meteorológica. Esta técnica estadístico-analítica parametriza las matrices de covarianza para tolerar patrones arbitrarios de correlación y heterocedasticidad dentro de cada clúster individual, derivando p-valores y límites de confianza robustos para la inferencia de tendencias.
    \subsubsection{Cálculo teórico de los errores estándar agrupados}

En el modelado de datos de panel, la presencia de correlación serial dentro de una misma estación meteorológica (autocorrelación) e idoneidades de varianza no constante (heterocedasticidad) invalida la inferencia basada en los errores estándar ordinarios, tendiendo a subestimar significativamente la verdadera variabilidad de los estimadores. Para corregir este problema, se implementa el estimador de matriz de covarianza robusta agrupada por clústeres (estaciones), el cual generaliza el estimador robusto de White para permitir correlación arbitraria entre las observaciones de una misma entidad \cite{wooldridge2010}.

Considérese el modelo lineal agrupado por estaciones (clústeres):
\begin{equation}
    y_{i} = X_{i}\beta + u_{i}, \quad i = 1, 2, \dots, M
\end{equation}
Donde $M$ representa el número total de clústeres (estaciones meteorológicas seleccionadas), $y_i$ es un vector de dimensión $T_i \times 1$ de la variable dependiente para la estación $i$, $X_i$ es la matriz de regresores de dimensión $T_i \times K$, y $u_i$ representa el vector de perturbaciones idiosincráticas de dimensión $T_i \times 1$.

El estimador de mínimos cuadrados ordinarios (o de efectos aleatorios tras la transformación de cuasi-diferenciación) se define como:
\begin{equation}
    \hat{\beta} = \left( \sum_{i=1}^{M} X_{i}' X_{i} \right)^{-1} \sum_{i=1}^{M} X_{i}' y_{i}
\end{equation}

La varianza asintótica de este estimador, condicionada en la matriz de regresores $X$, se modela bajo la estructura de varianza-covarianza tipo ``sándwich'':
\begin{equation}
    \text{Var}(\hat{\beta} | X) = \left( X'X \right)^{-1} B \left( X'X \right)^{-1}
\end{equation}
Donde el núcleo interno $B$ se define como la suma ponderada de las matrices de covarianza de los residuos para cada clúster independiente:
\begin{equation}
    B = \sum_{i=1}^{M} X_{i}' \Omega_{i} X_{i}
\end{equation}
En este esquema, $\Omega_i = E[u_i u_i' | X_i]$ es la matriz de varianza-covarianza de dimensión $T_i \times T_i$ específica para la estación $i$. El supuesto fundamental de este método es que las observaciones correspondientes a estaciones distintas son independientes entre sí ($E[u_i u_j' | X_i, X_j] = 0$ para todo $i \neq j$), permitiendo al mismo tiempo cualquier estructura de correlación y heterocedasticidad no paramétrica dentro de cada estación $i$.

Dado que $\Omega_i$ es desconocida, el estimador robusto agrupado sustituye esta matriz por el producto cruzado de los vectores de residuos muestrales ajustados para cada estación, $\hat{u}_i = y_i - X_i \hat{\beta}$:
\begin{equation}
    \widehat{\Omega}_{i} = \hat{u}_{i} \hat{u}_{i}'
\end{equation}

De este modo, la estimación empírica de la matriz de varianza-covarianza de los coeficientes, $\widehat{\text{Var}}_{robust}(\hat{\beta})$, se calcula mediante la fórmula:
\begin{equation}
    \widehat{\text{Var}}_{robust}(\hat{\beta}) = \left( \sum_{i=1}^{M} X_{i}' X_{i} \right)^{-1} \left( \sum_{i=1}^{M} X_{i}' \hat{u}_{i} \hat{u}_{i}' X_{i} \right) \left( \sum_{i=1}^{M} X_{i}' X_{i} \right)^{-1}
\end{equation}

Para corregir los sesgos inherentes en muestras donde el número de clústeres $M$ es finito, se aplica un factor de corrección de escala multiplicativo en la matriz de varianza-covarianza:
\begin{equation}
    q = \left( \frac{M}{M - 1} \right) \left( \frac{N - 1}{N - K} \right)
\end{equation}
Donde $N = \sum_{i=1}^{M} T_i$ representa el número total de observaciones mensuales del panel, y $K$ es el número de parámetros estimados en el modelo. Los errores estándar agrupados por estación corresponden a la raíz cuadrada de los elementos de la diagonal principal de la matriz $q \cdot \widehat{\text{Var}}_{robust}(\hat{\beta})$.
\end{enumerate}

Finalmente, una limitación inherente al marco de modelización propuesto reside en su especificación de tendencia lineal. El estimador asume una tasa de cambio constante en la anomalía hídrica a lo largo del periodo temporal estudiado, constituyendo una primera aproximación de tendencia a largo plazo que no captura dinámicas no lineales de escala decadal o cambios abruptos de régimen climático.

\chapter{Materiales y Métodos}

\section{Diseño de la investigación}

La presente investigación se inscribe en un enfoque cuantitativo, con un diseño no experimental, observacional, retrospectivo, longitudinal y explicativo. Es de carácter observacional y no experimental debido a que las variables climáticas se analizan de forma pasiva, tal como ocurrieron de manera natural en el entorno, sin manipulación deliberada. Posee un enfoque retrospectivo, puesto que se examinan registros históricos de un período extendido de 24 años (2000--2023). Asimismo, su diseño es longitudinal y explicativo mediante el uso de una estructura de datos de panel, orientada a estimar parámetros de tendencia de largo plazo para explicar y contrastar formalmente las hipótesis sobre las dinámicas de la sequía y el proceso de aridificación en Paraguay.

\section{Fuentes de datos y variables}

Los datos primarios corresponden a los registros meteorológicos diarios provistos por la Dirección de Meteorología e Hidrología (DMH) de la Dirección Nacional de Aeronáutica Civil (DINAC) del Paraguay. 

\subsection{Procedimiento de extracción y validación de Datos}

La serie temporal diaria analizada comprende el período desde el 1 de enero de 2000 hasta el 31 de diciembre de 2023. La fecha de extracción y recepción oficial de la base de datos fue el 9 de mayo de 2024 (Expediente N.º 213.042/2024; ver Apéndice \ref{apendice:fuente_datos}), tras una solicitud formal gestionada a través de la Facultad de Ciencias Exactas y Naturales (FACEN) de la Universidad Nacional de Asunción.

De forma complementaria, se recopilaron los metadatos de las estaciones terrestres, incluyendo su nombre oficial, identificador (ID Estación), coordenadas geográficas de posicionamiento (latitud y longitud decimal) y su macroregión geográfica correspondiente (Región Oriental o Región Occidental).

\subsection{Definición de variables climáticas}

El análisis de datos de panel se centra de manera exclusiva en las variables de precipitación diaria (lluvia) y temperatura media diaria como factores determinantes del balance hídrico. Las variables críticas recopiladas con frecuencia diaria para cada estación meteorológica fueron:
\begin{itemize}
    \item \textit{Precipitación diaria ($P_t$):} medida en milímetros ($mm$).
    \item \textit{Temperatura media diaria ($T_{\text{med}, t}$):} medida en grados Celsius ($^{\circ}\text{C}$).
    \item \textit{Temperatura máxima diaria ($T_{\text{máx}, t}$):} usada como variable de apoyo para la imputación térmica ($^{\circ}\text{C}$).
    \item \textit{Temperatura mínima diaria ($T_{\text{mín}, t}$):} usada como variable de apoyo para la imputación térmica ($^{\circ}\text{C}$).
\end{itemize}

\section{Procedimiento, depuración e imputación de datos}

Se examinó la cobertura histórica de cada una de las 19 estaciones terrestres respecto al rango teórico esperado de 8,766 días (24 años). El análisis de completitud determinó que la estación de Paraguarí (PARA) presentaba un truncamiento temporal crítico: sus registros inician únicamente a partir del 1 de abril de 2006, registrando una pérdida masiva superior al 26\% de la serie temporal bajo estudio (2,282 días faltantes). 

Con el fin de evitar un sesgo por truncamiento sistemático de la serie temporal nacional, se decidió excluir de manera definitiva la estación de Paraguarí (PARA). Las 18 estaciones restantes exhibieron tasas de completitud de fechas superiores al 97.9\%, conformando la muestra final del estudio. Sobre esta red de 18 estaciones se construyó un panel balanceado inicial de 157,788 registros diarios (18 estaciones $\times$ 8,766 días). Las filas correspondientes a fechas faltantes originalmente en cada estación fueron insertadas e inicializadas con valores nulos (\textit{NaN}).

\subsection{Algoritmo de imputación de temperatura por capas}

Para asegurar la continuidad temporal de la temperatura media diaria ($T {\text{med}}$) y evitar la pérdida de observaciones mensuales, se diseñó e implementó un algoritmo de imputación estructurado en tres capas jerárquicas:

\begin{enumerate}
    \item \textit{Capa 1: Imputación por relación térmica.} Para aquellos días donde se registró la ausencia de la temperatura media ($T {\text{med}}$) pero se disponía de los registros válidos de temperatura máxima ($T {\text{máx}}$) y mínima ($T {\text{mín}}$), se estimó el valor faltante como el promedio aritmético de las temperaturas extremas diarias:
    \begin{equation}
        T_{\text{med}, t} = \frac{T_{\text{mín}, t} + T_{\text{máx}, t}}{2}
    \end{equation}
    Este procedimiento permitió recuperar con precisión física un total de 342 días del panel.
    \item \textit{Capa 2: Interpolación lineal para vacíos cortos ($\le 3$ días).} Para los registros nulos persistentes de temperatura media, se analizó la longitud de los vacíos consecutivos. Para brechas de duración corta, definidas estrictamente entre 1 y 3 días consecutivos (que representan el 37.1\% del total de datos térmicos faltantes en la red), se aplicó una interpolación lineal temporal basada en la inercia térmica de los días adyacentes válidos. Esta capa recuperó un total de 655 días de registro.
    \item \textit{Capa 3: Imputación espacial por estación vecina ($\geq 4$ días).} Para vacíos de información térmica igual o mayor a cuatro días consecutivos, la interpolación temporal directa pierde validez física al ignorar la variabilidad sinóptica de mediano plazo. En consecuencia, se aplicó un método de transferencia espacial basado en la cercanía geográfica. 

Utilizando las coordenadas geográficas de las 18 estaciones meteorológicas seleccionadas, se calculó una matriz de distancias mutuas (ortodrómicas) que permitió identificar de forma unívoca a la estación vecina más cercana para cada punto de la red (por ejemplo, la menor distancia registrada correspondió al par Coronel Oviedo-Villarrica con solo $25.75\text{ km}$, mientras que la mayor distancia de vecindad se registró en el Chaco Profundo entre Mariscal Estigarribia y Pozo Colorado con $248.84\text{ km}$). 

Para cada fecha que presentaba vacíos térmicos prolongados en la estación objetiva, se transfirió directamente el registro diario de temperatura media depurada ($T_{\text{med\_imputada}}$) correspondiente a su vecina más cercana en ese mismo día. Este procedimiento de interpolación espacial permitió recuperar con precisión climatológica un total de 1,114 observaciones diarias. Al término de esta tercera capa jerárquica de imputación, se logró una tasa de completitud global en las series de temperatura media diaria del 99.99\%, registrando únicamente 9 datos nulos residuales en toda la base de datos de panel maestro para las 18 estaciones a lo largo de los 24 años de estudio.
\end{enumerate}

\subsection{Criterio de no Imputación para la variable precipitación}

A diferencia de la temperatura, la precipitación diaria es una variable altamente asimétrica, intermitente y no continua. Debido a que las fallas instrumentales y los vacíos de transmisión de la red meteorológica ocurren por factores técnicos o administrativos totalmente independientes del estado del tiempo, no existe relación causal entre la omisión del dato y el evento físico de lluvia. Bajo esta premisa, como lo afirma \citeA{wooldridge2010}, el sesgo no depende de la cantidad o el patrón de los datos faltantes, sino de la naturaleza del mecanismo de selección que genera esos nulos. Si la falta de datos es aleatoria o técnica, el estimador de efectos fijos manejará correctamente el desbalance del panel sin introducir sesgos de truncamiento temporal.

Consecuentemente, se adoptó la decisión metodológica de no aplicar ningún método de imputación sobre la variable precipitación. Forzar estimaciones (ya sean interpolaciones lineales o vecindades) sobre la lluvia introduciría sesgos artificiales severos, alterando la frecuencia real de días lluviosos y la magnitud acumulada de los eventos.

\section{Agregación de datos y creación de la base mensual}

Dado que el balance hídrico y el cálculo del SPEI se analizan a nivel mensual, el panel diario depurado de 18 estaciones fue agregado a una frecuencia mensual de 288 meses consecutivos por estación (5,184 observaciones mensuales potenciales). 

Para garantizar la fiabilidad física de los agregados mensuales, se implementó la regla de exclusión recomendada por la Organización Meteorológica Mundial \cite{wmo2017}:
\begin{itemize}
    \item \textit{lluvia\_total:} suma mensual de la precipitación diaria. Se define como nula (\textit{NaN}) si el mes correspondiente registra menos de 25 días con observaciones de lluvia completas.
    \item \textit{temp\_media:} promedio mensual de la temperatura media diaria imputada. Se define como nula si el mes registra menos de 25 días con datos térmicos válidos.
\end{itemize}

Este procedimiento determinó un panel mensual final caracterizado por una alta completitud: la serie de temperatura media mensual reportó un único dato nulo ($0.01\%$), mientras que la serie de precipitación mensual acumulada presentó únicamente 82 registros nulos ($1.58\%$) de un total de 5,184 observaciones.

\section{Construcción de la variable dependiente: el índice SPEI-6}

Para el cálculo del Índice de Precipitación-Evapotranspiración Estandarizado a escala de 6 meses (SPEI-6) sobre las 18 estaciones finales, se ejecutó secuencialmente la siguiente formulación matemática:

\begin{enumerate}
    \item \textit{Cálculo de la ETP mensual (Thornthwaite):} Estimación de la demanda evaporativa potencial mediante el método empírico de \citeA{thornthwaite_approach_1948}, utilizando la temperatura media mensual ($T_{\text{avg}}$) y la latitud geográfica de cada estación para derivar el exponente empírico $a$ y el factor de corrección latitudinal mensual $K$.
    \item \textit{Cálculo del balance hídrico mensual ($D_i$):} Definido como la diferencia simple entre la precipitación total mensual ($P_i$) y la ETP ajustada ($ETP_i$) para cada mes $i$: $D_i = P_i - ETP_i$.
    \item \textit{Acumulación multiescalar ($D^k_i$):} Para capturar los efectos de mediano plazo se acumuló el balance hídrico mensual en una ventana temporal de $k=6$ meses de forma consecutiva \cite{mckee1993}.
    \item \textit{Ajuste paramétrico Pearson Tipo III:} Los valores acumulados de balance hídrico ($D^6_i$) fueron ajustados a una distribución teórica Pearson Tipo III (Gamma de tres parámetros) para estimar la probabilidad acumulada $P = F(x)$, lo cual resulta óptimo para tolerar asimetrías y valores negativos de balance.
    \item \textit{Estandarización Normal estándar:} La probabilidad acumulada $P$ se transformó en anomalías estandarizadas normales ($Z$-scores) mediante la aproximación racional analítica de \citeA{abramowitz_handbook_1965}. El valor final resultante constituye la variable dependiente de panel, SPEI-6.
\end{enumerate}

\section{Regionalización geográfica mediante agrupamiento no supervisado}

Para analizar la distribución espacial de la sequía y evitar límites políticos arbitrarios, se aplicó un algoritmo de agrupamiento jerárquico aglomerativo (HAC) sobre las series de SPEI-6 de las 18 estaciones. Se implementó la distancia de disimilitud euclídea y el criterio de enlace de varianza mínima de Ward. El número óptimo de subregiones climáticas ($k^*$) se determinó mediante la maximización del Coeficiente de Silueta Promedio de Rousseeuw \cite{rousseeuw1987}.

\section{Especificación del modelo de panel y test de Hausman}\label{sec:modelo_estadistico}

Para cuantificar las tendencias temporales del balance hídrico se ajustó el modelo de regresión lineal de panel especificado en la ecuación \ref{eq:modelo_completo} y la selección entre la estimación por efectos fijos (FE) y efectos aleatorios (RE) se fundamentó en el test de especificación de \citeA{hausman1978}, evaluando formalmente la consistencia del estimador RE frente a la ortogonalidad de los efectos individuales no observados de las 18 estaciones meteorológicas. Para esto se utilizó la librería \textit{linearmodels} de Python.

La especificación completa del modelo estimado, incluyendo todas las variables e interacciones, fue la siguiente:
\sloppy
\begin{equation}
\begin{split}
    \text{SPEI-6}_{it} = & \beta_0 + \beta_1 \text{tiempo}_t + \sum_{s=1}^{3} \gamma_s \text{estacion-climatica}_{s,t} + \delta_1 \text{region-geografica}_i \\
    & + \sum_{s=1}^{3} \theta_s (\text{tiempo}_t \times \text{estacion-climatica}_{s,t}) + \phi_1 (\text{tiempo}_t \times \text{region-geografica}_i) \\
    & + \sum_{s=1}^{3} \psi_s (\text{estacion-climatica}_{s,t} \times \text{region-geografica}_i) \\
    & + \sum_{s=1}^{3} \omega_s (\text{tiempo}_t \times \text{estacion-climatica}_{s,t} \times \text{region-geografica}_i) + u_{it}
\end{split}
\label{eq:modelo_completo}
\end{equation}
donde $u_{it} = \alpha_i + \varepsilon_{it}$ es el término de error compuesto, y los términos $\text{estacion-climatica}_{s,t}$ y $\text{region-geografica}_i$ son variables dummys para las categorías estacionales climáticas (invierno, primavera, verano, otoño) y regionales geográficas (clusters), respectivamente.

\subsection{Corrección de la dependencia espacial: el estimador de Driscoll y Kraay}

Para controlar y corregir de manera simultánea la presencia de autocorrelación serial temporal y de correlación espacial o de sección cruzada identificada en la red de monitoreo, se prescinde de los errores estándar de clúster convencionales y se adopta el estimador de matriz de covarianza robusto desarrollado por \cite{driscoll1998}. 

El método de Driscoll-Kraay es un estimador no paramétrico de matriz de covarianza consistente ante la presencia de heterocedasticidad, autocorrelación y cualquier forma de dependencia espacial o de sección cruzada. Su ventaja metodológica radica en que estima la estructura de covarianza de las medias temporales de los momentos transversales mediante un estimador tipo Newey-West con un kernel espectral de Bartlett, asegurando la consistencia de los errores estándar y garantizando intervalos de confianza robustos para la inferencia de las tendencias dinámicas de sequía a nivel nacional \cite{wooldridge2010}.

\section{Marco de contraste SPEI vs. SPI para el aislamiento del efecto térmico}

Para dar cumplimiento a nuestro objetivo específico y contrastar de manera formal la hipótesis 3 (H3), se diseñó un marco experimental comparativo. El procedimiento consiste en estimar de forma paralela la ecuación del modelo de regresión de panel de efectos aleatorios (Ecuación \ref{eq:modelo_completo}) sobre dos variables dependientes distintas:
\begin{align}
    \text{SPEI-6}_{it} &= \mathbf{x}_{it}'\boldsymbol{\beta}_{SPEI} + u_{it} \label{eq:panel_spei} \\
    \text{SPI-6}_{it} &= \mathbf{x}_{it}'\boldsymbol{\beta}_{SPI} + v_{it} \label{eq:panel_spi}
\end{align}

Donde $\boldsymbol{\beta}_{SPEI}$ y $\boldsymbol{\beta}_{SPI}$ representan los vectores de parámetros estimados para los modelos de balance hídrico y precipitación univariada, respectivamente. El efecto aislado de la ETP sobre la dinámica temporal del balance hídrico se define para cada estación del año y región como la diferencia entre sus respectivas tendencias netas mensuales calculadas ($\Delta \beta_{neto} = \beta_{neto, SPEI} - \beta_{neto, SPI}$). 

La significancia estadística del efecto del aumento de la demanda evaporativa impulsada por el aumento de las temperaturas (calentamiento global) ($H_0: \beta_{neto, SPEI} - \beta_{neto, SPI} = 0$) se evalúa mediante la estimación de intervalos de confianza del 95\% y pruebas de Wald para diferencias de coeficientes de pendiente comunes entre ambas regresiones de panel, aislando así el efecto puro de la temperatura por encima del régimen de precipitaciones.

\chapter{RESULTADOS Y DISCUSIÓN}

\section{Regionalización climatológica de la sequía en Paraguay}

Antes de proceder a caracterizar los eventos de sequía y estimar tendencias dinámicas mediante el modelado de datos de panel, examinamos desde una perspectiva empírica y de balance hídrico la variabilidad espacial de la sequía en Paraguay.

Con el propósito de establecer subregiones climáticas homogéneas de manera objetiva, se aplicó un algoritmo de agrupamiento jerárquico aglomerativo sobre la matriz de datos continuos mensuales del índice SPEI-6 para las 18 bases meteorológicas seleccionadas durante el período 2000--2023. El preprocesamiento de la base detectó 90 celdas con valores perdidos en las series temporales originales del SPEI-6, los cuales fueron resueltos mediante interpolación lineal temporal antes de proceder con el agrupamiento.

\subsection{Optimización y selección del número de clústeres}

La jerarquía de fusiones consecutivas calculada por el algoritmo de Ward se representa de manera gráfica en el dendrograma de disimilitud de la Figura \ref{fig:regionalizacion} (Panel A). Para evitar la discrecionalidad en el corte jerárquico de este dendrograma, se calculó el coeficiente de silueta promedio para un rango operativo de $k \in [2, 5]$ clústeres. La Tabla \ref{tab:silueta} expone los resultados de esta optimización espacial, la cual se ilustra además en la curva de optimización de la Figura \ref{fig:regionalizacion} (Panel B).

El coeficiente de silueta promedio alcanza su máximo global para una estructura de $k^* = 2$ clústeres ($s = 0.1758$), indicando que una partición de dos regiones climáticas proporciona la máxima cohesión intra-grupo y la mayor diferenciación inter-grupo. Las particiones alternativas de $3$, $4$ o $5$ grupos muestran un decaimiento progresivo en la calidad de la silueta, justificando estadísticamente la selección de dos subregiones óptimas a nivel nacional para el período bajo estudio.

\begin{table}[ht!]
\centering
\caption{Coeficiente de silueta promedio según el número de particiones (clústeres).}
\label{tab:silueta}
\vspace{-0.6cm}
\singlespacing % <-- ACTIVA INTERLINEADO SENCILLO LOCAL
\small         % <-- ACTIVA TAMAÑO DE LETRA CHICA LOCAL (10pt)
\begin{tabular}{cc}
\hline\hline
\textbf{Número de clústeres ($k$)} & \textbf{Coeficiente de silueta promedio} \\ \hline
$\mathbf{k = 2}$ & \textbf{0.1758} \\
$k = 3$ & 0.1682 \\
$k = 4$ & 0.1548 \\
$k = 5$ & 0.1592 \\ \hline\hline
\end{tabular}
\end{table}

\begin{figure}[htbp]
\centering
\includegraphics[width=\textwidth]{/home/hjvargaso/proyectos/tesis_sequia_2/output/figures/analisis_regionalizacion_sequia.png}
\vspace{-0.75cm}
\begin{minipage}{14cm}
\caption{Resultados de la regionalización climática de la sequía en Paraguay (periodo 2000--2023): (A) Dendrograma de agrupamiento jerárquico aglomerativo utilizando la distancia de disimilitud de Ward; (B) Optimización del número de clústeres ($k$) mediante el coeficiente de silueta promedio; (C) Distribución espacial resultante de las bases meteorológicas terrestres clasificadas por subregión climática.}
\vspace{0.5cm}
\label{fig:regionalizacion}
\end{minipage}
\end{figure}

\subsection{Asignación de bases meteorológicas y validación de contigüidad geográfica}

Para validar la coherencia climática y física de la regionalización propuesta, se contrastó la asignación temporal derivada del clustering de SPEI-6 con la ubicación geográfica (latitud y longitud) de las bases terrestres. La evaluación cuantitativa de las distancias de posicionamiento reveló que la distancia geográfica promedio entre bases asignadas al mismo clúster (distancia intra-clúster) es de $2.438^{\circ}$ decimales (aproximadamente $260$ km), mientras que la distancia promedio entre bases pertenecientes a clústeres distintos (distancia inter-clúster) asciende a los $3.322^{\circ}$ decimales (aproximadamente $360$ km).

Esta diferencia confirma la existencia de una contigüidad geográfica compacta. Como se presenta visualmente en el mapa de la Figura \ref{fig:regionalizacion} (Panel C), la clasificación basada en el SPEI-6 ha particionado el territorio nacional de manera limpia, delimitando dos subregiones climáticas claramente definidas y coherentes con sus características de balance hídricas:

\begin{enumerate}
    \item \textit{Subregión sureste (Clúster 1 - naranja):} Integrada por las bases de Caazapá (CAAZ), Villarrica (VILL), Coronel Oviedo (COVIE), San Juan Bautista (SJBA), Minga Guazú (MINGA), Capitán Meza (CMEZA) y Encarnación (ENCAR). Este bloque representa la zona sur y este de la región oriental, históricamente caracterizada por ser el área de mayor pluviometría de Paraguay, con balances hídricos promedio superavitarios y una menor inercia a eventos de sequía extrema prolongada.
    \item \textit{Subregión centrochaco (Clúster 2 - verde):} Conformada por las bases de Mariscal Estigarribia (MCAL), Puerto Casado (PTOC), Pozo Colorado (PCOL), Gral. José María Bruguez (GBRU), Concepción (CONC), San Pedro (SPED), San Estanislao (SEST), Salto del Guairá (SGUA), Luque (LUQUE), Pilar (PILAR) y Pedro Juan Caballero (PJC). Esta amplia franja geográfica continua se extiende desde el chaco profundo y semiárido, cruzando el río Paraguay hacia el centronorte y norte de la región oriental, incluyendo además bases de transición en el sur de la cuenca (Pilar). Estas bases comparten una estación seca invernal más marcada y promedios térmicos elevados que incrementan sistemáticamente la demanda evaporativa.
\end{enumerate}

Esta partición de dos subregiones climáticas óptimas reemplazará la división regional clásica (región oriental y occidental) en todos los análisis descriptivos e inferenciales subsiguientes de la investigación, permitiendo examinar el proceso de aridificación con un mayor grado de precisión espacial y de balance hídrico.

\section{Patrones estacionales del balance hídrico}

El análisis de la distribución mensual del SPEI-6, presentado mediante diagramas de caja en la Figura \ref{fig:distribucion_mensual}, revela la existencia de un ciclo estacional marcado en las condiciones hídricas del país. La variabilidad observada a lo largo de los meses muestra que Paraguay experimenta transiciones significativas entre períodos de superávit y de déficit hídrico acumulado.

\begin{figure}[htbp]
    \centering
    \includegraphics[width=0.85\textwidth]{/home/hjvargaso/proyectos/tesis_sequia_2/output/figures/distribucion_mensual_spei6.png}
    \vspace{-0.25cm}
    \begin{minipage}{14cm}
    \centering
    \caption{Distribución mensual del índice SPEI-6 para las 18 bases seleccionadas durante el período 2000--2023. La línea horizontal negra dentro de cada caja representa la mediana de la distribución; los límites naranja y rojo indican los umbrales de sequía moderada ($\leq -1.0$) y severa ($\leq -1.5$), respectivamente. Los valores atípicos se grafican como puntos grises fuera de los bigotes.}
    \label{fig:distribucion_mensual}
    \end{minipage}
\end{figure}

Como se ilustra en la Figura \ref{fig:distribucion_mensual}, el final del invierno y el inicio de la primavera (agosto y septiembre) constituyen climatológicamente los meses más secos a nivel de base. Durante este bimestre, las medianas del SPEI-6 se posicionan de manera sistemática en valores negativos, siendo las únicas del ciclo anual situadas por debajo de la línea de neutralidad hídrica (SPEI = 0). Este comportamiento físico se corresponde con la marcada disminución estacional de las precipitaciones en el territorio nacional durante el período frío, lo que reduce la humedad disponible en el suelo.

En contraste, los meses correspondientes al final del verano y el otoño (de febrero a mayo) se perfilan consistentemente como los períodos con mayor superávit hídrico, exhibiendo medianas del SPEI-6 con desviaciones positivas. No obstante, un detalle de gran relevancia para la caracterización de las sequías en Paraguay radica en la distribución de los valores atípicos (extremos) durante los meses de verano, específicamente en diciembre y enero. 

A pesar de que el verano se caracteriza en promedio por un balance hídrico neutral a ligeramente húmedo, los valores mínimos extremos del índice —aquellos que superan el umbral de sequía extrema (SPEI-6 $\leq -2.0$) y alcanzan anomalías severas inferiores a $-3.5$— se registran con frecuencia durante el período estival. Este comportamiento evidencia que, aunque las lluvias de verano suelen ser cuantitativamente abundantes en Paraguay, las anomalías térmicas positivas y las olas de calor que ocurren de manera sincrónica en esta estación pueden disparar de tal forma la demanda evaporativa de la atmósfera (ETP) que superan el volumen de precipitaciones. De este modo, se configuran sequías estivales rápidas y severas que actúan como propulsoras del proceso de aridificación estacional.

\section{Caracterización descriptiva de las dinámicas de sequía (2000--2023)}

Para comprender las características de las sequías en Paraguay, esta sección aborda cinco dimensiones analíticas fundamentales: estacionalidad anual, frecuencia de eventos, intensidad, duración temporal y extensión espacial de las anomalías de balance hídrico.  Los resultados estadísticos se sintetizan de forma gráfica en la Figura \ref{fig:caracterizacion}.

\begin{figure}[htbp]
    \centering
    \begin{minipage}{14cm}
    \includegraphics[width=\textwidth]{/home/hjvargaso/proyectos/tesis_sequia_2/output/figures/caracterizacion_descriptiva_sequia.png}
    \caption{Caracterización de las dinámicas de sequía en Paraguay (2000--2023): (A) Frecuencia anual de meses-bases en sequía moderada y severa; (B) Evolución de la intensidad promedio y mínima anual; (C) Distribución comparativa de la duración de eventos entre el Periodo 1 (2000--2011) y el Periodo 2 (2012--2023); (D) Extensión espacial promedio anual (\% de bases afectadas simultáneamente).}
    \label{fig:caracterizacion}
    \end{minipage}
\end{figure}

\subsection{Análisis de frecuencia de eventos extremos}

Para cuantificar los cambios en la ocurrencia de la sequía, se evaluó la frecuencia anual de ``meses-bases'' que superaron los umbrales de sequía moderada ($SPEI-6 \leq -1.0$) y sequía severa/extrema ($SPEI-6 \leq -1.5$). Como se ilustra en la Figura \ref{fig:caracterizacion} (Panel A), existe una tendencia positiva y sostenida en la frecuencia de estos eventos a lo largo de los 24 años analizados.

Durante la primera década del estudio (2000--2009), la ocurrencia de meses en sequía se mantuvo dentro de los límites de variabilidad interanual normal del país. No obstante, a partir de 2019 se observa un quiebre estructural sin precedentes en la serie histórica. El pico máximo de frecuencia se registró en el año 2020, acumulando una cifra récord de 80 meses-bases bajo condiciones de sequía moderada o severa. Considerando que el total anual posible de meses-bases para la red de 18 bases es de 216 ($18 \text{ bases} \times 12 \text{ meses}$), este resultado indica que aproximadamente el 37\% del territorio nacional monitoreado experimentó déficit hídrico de forma simultánea a lo largo del año 2020, consolidando la crisis hídrica más generalizada de las últimas dos décadas.

\subsection{Análisis de intensidad del déficit hídrico}

La intensidad de la sequía se examinó mediante dos descriptores anuales: la intensidad promedio de los meses en sequía ($SPEI-6 \leq -1.0$) y el valor mínimo absoluto registrado en la red de bases. Como se expone en la Figura \ref{fig:caracterizacion} (Panel B), ambas métricas muestran una trayectoria de agravamiento del déficit.

El promedio anual de la intensidad de las sequías ha transitado de valores cercanos a $-1.15$ a inicios del período estudiado, hacia promedios sistemáticamente inferiores a $-1.60$ en los últimos años del registro. El mínimo absoluto de toda la serie histórica se registró en diciembre de 2023 (2023-12), alcanzando un valor extremo de $-3.7103$ en la base de Mariscal Estigarribia (MCAL), ubicada en el chaco central. Este registro, que se sitúa a casi cuatro desviaciones estándar por debajo de la normal climatológica, supera ampliamente el umbral técnico de sequía extrema ($SPEI-6 \leq -2.0$). Este comportamiento ilustra la severidad que adquieren las sequías de Paraguay cuando las anomalías de temperatura potencian la demanda evaporativa en ausencia de precipitaciones, actuando como un impulsor físico del proceso de aridificación local.

\subsection{Análisis de la duración temporal (persistencia)}

La duración de los eventos de sequía (persistencia temporal) se evaluó mediante un análisis de rachas, definiendo la duración de un evento como el número de meses consecutivos en que una base meteorológica registra un valor de $SPEI-6 \leq -1.0$. Bajo este esquema metodológico, la duración promedio general de los eventos de sequía para el territorio nacional durante los 24 años analizados se ubicó en $2.96$ meses, registrándose una racha máxima absoluta de $14$ meses continuos de déficit hídrico extremo. En la Figura \ref{fig:caracterizacion} (Panel C) se presenta la distribución de estas duraciones, acotada visualmente a un rango de 12 meses para preservar la escala de los histogramas.

Para examinar si las fluctuaciones climáticas de las últimas décadas han alterado la persistencia temporal de estos fenómenos, se contrastaron los resultados agregados del primer intervalo de estudio (Periodo 1: 2000--2011) frente al intervalo de calentamiento más pronunciado (Periodo 2: 2012--2023). 

El análisis comparativo de la distribución de frecuencias de las rachas de sequía evidencia un cambio estructural notable. Durante el Periodo 1, la duración promedio de las sequías fue de $2.80$ meses, caracterizándose por una alta frecuencia de eventos secos transitorios de corta duración (el histograma del Panel C exhibe un pico marcado de más de 55 eventos de un solo mes de duración). En contraste, para el Periodo 2, la duración promedio se incrementó a $3.12$ meses (un aumento relativo del 12\% en la persistencia del fenómeno), registrándose una disminución en las sequías de un mes y un incremento sistemático en los eventos persistentes de 2 a 4 meses de duración consecutiva.

Esta prolongación de las rachas consecutivas indica que, en las últimas dos décadas, los sistemas de balance hídrico terrestre en Paraguay han experimentado una pérdida en su capacidad de recuperación rápida. Esto prolonga los estados de estrés vegetativo y agota las reservas de humedad del suelo de manera acumulativa, lo que constituye un claro indicador físico del proceso de aridificación que experimenta la región.

\subsection{Análisis de la extensión espacial de la sequía}

La extensión espacial se cuantificó como el porcentaje promedio de bases de la red que experimentaron condiciones de sequía simultáneamente ($SPEI-6 \leq -1.0$) en cada año calendario. El promedio anual de afectación y su correspondiente línea de tendencia lineal se presentan en la Figura \ref{fig:caracterizacion} (Panel D).

Los resultados confirman que la sequía en Paraguay ha transitado de ser un fenómeno eminentemente zonal o subregional a principios del milenio, a convertirse en una anomalía de escala nacional y más generalizada durante la última década. La línea de tendencia muestra una pendiente positiva pronunciada que inicia cerca del 10\% de afectación espacial en el año 2000 y se proyecta por encima del 20\% para el año 2023. Durante el ciclo crítico de 2020--2021, la extensión espacial del déficit hídrico alcanzó niveles históricos sin precedentes, afectando de forma simultánea a un promedio anual superior al 35\% de la red de monitoreo nacional.

\section{Resultados del modelo de datos de panel}

Para cuantificar y evaluar las tendencias temporales del balance hídrico en Paraguay, se estimó el modelo de regresión lineal para datos de panel con interacciones múltiples especificado en la Ecuación \ref{eq:modelocompleto}. Las series temporales de las 18 bases meteorológicas seleccionadas cubren un período de $288$ meses consecutivos (2000--2023), totalizando un tamaño muestral balanceado de $5184$ observaciones mensuales.

\subsection{Selección del estimador: El test de especificación de Hausman}

La elección metodológica entre el modelo de regresión de efectos fijos (FE) y el de efectos aleatorios (RE) constituye una decisión crítica para garantizar la consistencia y la eficiencia de los parámetros estimados \cite{wooldridge2010}. Los resultados numéricos del test de Hausman se detallan en la Tabla \ref{tab:test_hausman}.

\begin{table}[ht!]
\centering
\caption{Resultados del test de especificación de Hausman para la selección del estimador.}
\label{tab:test_hausman}
\vspace{-0.6cm}
\singlespacing % <-- ACTIVA INTERLINEADO SENCILLO LOCAL
\small         % <-- ACTIVA TAMAÑO DE LETRA CHICA LOCAL (10pt)
\begin{tabular}{p{6cm}c}
\hline\hline
\textbf{Test} & \textbf{Valor} \\ \hline
Estadístico de Prueba ($\chi^2$) & -0.0035 \\
Grados de Libertad ($k$) & 15 \\
p-valor & 1.0000 \\ \hline\hline
\end{tabular}
\end{table}

De acuerdo con los resultados expuestos en la Tabla \ref{tab:test_hausman}, el estadístico de prueba de Hausman arroja un valor de $-0.0035$, con un p-valor correspondiente de $1.0000$. En la práctica econométrica con muestras finitas, la obtención de un estadístico de Hausman negativo es un fenómeno común que ocurre cuando la diferencia entre las matrices de covarianza de ambos modelos ($\boldsymbol{V}_{FE} - \boldsymbol{V}_{RE}$) no es definida positiva \cite{baltagi2021}. Conforme con la literatura de referencia \cite{hausman1981}, este resultado se interpreta formalmente como la imposibilidad de rechazar la hipótesis nula de ortogonalidad, convalidando la consistencia del estimador de efectos aleatorios.

Un p-valor que converge de forma exacta a la unidad ($1.0000$) indica que la discrepancia entre el vector de coeficientes de efectos fijos ($\boldsymbol{\hat{\beta}}_{FE}$) y el de efectos aleatorios ($\boldsymbol{\hat{\beta}}_{RE}$) es infinitesimal. Este comportamiento responde a la naturaleza física y estadística de las variables implementadas en esta investigación:

\begin{enumerate}
    \item \textit{Efecto de la estandarización del SPEI-6:} Dado que la variable dependiente de la investigación ($y_{it} = \text{SPEI-6}_{it}$) es un índice climatológico estandarizado ($Z$-score), su nivel de base local o climatología media fue removida matemáticamente durante el proceso de ajuste paramétrico original \cite{vicenteserrano2010}. Al haberse neutralizado previamente la media climatológica específica de cada una de las 18 baese meteorológicas, la varianza residual asociada a los efectos individuales no observados ($\alpha_i$) tiende a comportarse de forma estrictamente ortogonal con respecto a los regresores explicativos de tendencia y estacionalidad.
    \item \textit{Consistencia y eficiencia de los coeficientes de pendiente:} La ortogonalidad demostrada garantiza que los coeficientes del modelo de efectos aleatorios no se encuentran sesgados por variables omitidas de carácter espacial invariante. En consecuencia, la aproximación de efectos aleatorios (RE) constituye el estimador óptimo para esta investigación. Al explotar simultáneamente la variabilidad temporal intra-base (\textit{within}) e inter-base (\textit{between}), el estimador RE proporciona una eficiencia asintótica superior para inferir de manera robusta las pendientes de aridificación estacional y regional en el Paraguay.
\end{enumerate}

\subsection{Verificación de la dependencia de sección cruzada (correlación espacial)}

Para verificar empíricamente la presencia de este fenómeno, se extrajeron los residuos de la estimación preliminar del modelo de efectos aleatorios y se aplicó la prueba de Dependencia de Sección Cruzada (CD) de Pesaran. El test evalúa la hipótesis nula ($H_0$) de independencia espacial o ausencia de correlación de sección cruzada entre las 18 bases meteorológicas del panel. 

La Tabla \ref{tab:test_pesaran} resume los resultados del test de Pesaran, rechazando la hipótes nula de independencia cruzada. Este resultado demuestra la existencia de una fuerte correlación espacial/sincrónica entre los residuos de la red de bases meteorológicas, lo cual exige metodológicamente la aplicación de un estimador de covarianza robusto como el de \citeA{driscoll1998} para garantizar la validez de la inferencia de las tendencias de largo plazo.

\begin{table}[ht!]
\centering
\begin{threeparttable}
\caption{Resultados de la prueba de dependencia de sección cruzada (CD) de Pesaran (2004) sobre los residuos del panel.}
\label{tab:test_pesaran}
\vspace{-0.6cm}
\singlespacing % <-- ACTIVA INTERLINEADO SENCILLO LOCAL
\small         % <-- ACTIVA TAMAÑO DE LETRA CHICA LOCAL (10pt)
\begin{tabular}{lcc}
\hline\hline
\textbf{Prueba Estadística} & \textbf{Estadístico CD} & \textbf{p-valor} \\ \hline
Dependencia de Sección Cruzada de Pesaran & 70.5878 & 0.0000 \\ \hline
\textbf{Hipótesis Nula ($H_0$)} & \multicolumn{2}{p{6cm}}{Independencia de sección cruzada (ausencia de correlación espacial).} \\
\textbf{Decisión Estadística} & \multicolumn{2}{p{6cm}}{Rechazar $H_0$ a un nivel de significancia de $\alpha = 0.01$.} \\
\textbf{Implicación Metodológica} & \multicolumn{2}{p{6cm}}{Obligatoriedad de implementar errores estándar robustos de Driscoll-Kraay (1998).} \\ \hline\hline
\end{tabular}
\end{threeparttable}
\end{table}

\subsection{Estimación del modelo final}

Una vez confirmada la idoneidad del estimador de efectos aleatorios mediante el test de Hausman, se procedió a estimar la especificación completa del modelo de regresión (Ecuación \ref{eq:modelocompleto}). Para corregir de forma preventiva la presencia de autocorrelación serial intragrupo y heterocedasticidad, y controlando la fuerte dependencia espacial contemporánea demostrada por el test de \citeA{pesaran2004}, el modelo se estimó utilizando errores estándar de \citeA{driscoll1998} robustos a la correlación espacial. La Tabla \ref{tab:estadisticas_modelo} resume las estadísticas generales de ajuste de la regresión.

\begin{table}[ht!]
\centering
\caption{Estadísticas de ajuste del modelo de efectos aleatorios estimado con errores robustos de Driscoll y Kraay (1998).}
\label{tab:estadisticas_modelo}
\vspace{-0.6cm}
\singlespacing % <-- ACTIVA INTERLINEADO SENCILLO LOCAL
\small         % <-- ACTIVA TAMAÑO DE LETRA CHICA LOCAL (10pt)
\begin{tabular}{p{8cm}c}
\hline\hline
\textbf{Métrica de Ajuste} & \textbf{Valor del Estimador} \\ \hline
Número de Observaciones ($N \times T$) & 5094 \\
R-cuadrado global ($R^2$ overall) & 0.0356 \\
R-cuadrado dentro ($R^2$ within) & 0.0356 \\
R-cuadrado entre ($R^2$ between) & 0.1903 \\
Estadístico F conjunto robusto & 1.9807 \\
p-valor de la regresión conjunta & 0.0132 \\
\hline\hline
\end{tabular}
\end{table}

El modelo en su conjunto es estadísticamente significativo, con un p-valor del estadístico F conjunto robusto de $0.0132$. Aunque el coeficiente de determinación global y dentro es modesto ($R^2 = 0.0356$), este comportamiento es consistente y esperado en modelos climáticos basados en anomalías de alta frecuencia, donde el componente de variabilidad a corto plazo introduce un volumen significativo de ruido blanco que no disminuye la consistencia ni la validez de los estimadores de pendiente de largo plazo \cite{stocker_climate_2013}. No obstante, cabe destacar el poder explicativo inter-base ($R^2 \text{ between} = 0.1903$), el cual denota que el modelo captura de forma consistente las diferencias sistemáticas en el balance hídrico promedio entre los clústeres geográficos de Paraguay.

La Tabla \ref{tab:resultados_re} expone los coeficientes estimados para cada regresor, sus errores estándar y su significancia estadística.

\begin{table}[ht!]
\centering
\caption{Resultados del modelo de panel con efectos aleatorios para la variable dependiente SPEI-6. Errores estándar de Driscoll y Kraay (1998) robustos a la correlación espacial.}
\label{tab:resultados_re}
\vspace{-0.6cm}
\singlespacing % <-- ACTIVA INTERLINEADO SENCILLO LOCAL
\small         % <-- ACTIVA TAMAÑO DE LETRA CHICA LOCAL (10pt)
\resizebox{\textwidth}{!}{%
\begin{tabular}{lcccc}
\hline\hline
\textbf{Regresor} & \textbf{Coeficiente} & \textbf{Error Est.} & \textbf{Estadístico t} & \textbf{p-valor} \\ \hline
intercepto ($\beta_0$) & -0.1844 & 0.1513 & -1.22 & 0.2232 \\
otoño ($\gamma_1$) & 0.3050 & 0.1892 & 1.61 & 0.1070 \\
primavera ($\gamma_2$) & 0.0892 & 0.2071 & 0.43 & 0.6665 \\
verano ($\gamma_3$) & 0.3447 & 0.2265 & 1.52 & 0.1282 \\
clúster 1 ($\delta_1$) & -0.0640 & 0.2174 & -0.29 & 0.7684 \\
otoño $\times$ clúster 1 ($\psi_1$) & 0.0158 & 0.2522 & 0.06 & 0.9502 \\
primavera $\times$ clúster 1 ($\psi_2$) & 0.3467 & 0.2860 & 1.21 & 0.2254 \\
verano $\times$ clúster 1 ($\psi_3$) & 0.4448 & 0.3620 & 1.23 & 0.2191 \\
Tiempo ($\beta_1$) & 0.0012 & 0.0012 & 1.00 & 0.3150 \\
Tiempo $\times$ otoño ($\theta_1$) & -0.0002 & 0.0015 & -0.15 & 0.8833 \\
Tiempo $\times$ primavera ($\theta_2$) & -0.0023 & 0.0015 & -1.54 & 0.1244 \\
Tiempo $\times$ verano ($\theta_3$) & -0.0022 & 0.0017 & -1.27 & 0.2056 \\
Tiempo $\times$ clúster 1 ($\phi_1$) & -0.0002 & 0.0017 & -0.12 & 0.9008 \\
Tiempo $\times$ otoño $\times$ clúster 1 ($\omega_1$) & -0.0015 & 0.0017 & -0.86 & 0.3890 \\
Tiempo $\times$ primavera $\times$ clúster 1 ($\omega_2$) & 0.0003 & 0.0023 & 0.15 & 0.8846 \\
Tiempo $\times$ verano $\times$ clúster 1 ($\omega_3$) & -0.0018 & 0.0025 & -0.74 & 0.4587 \\ \hline\hline
\multicolumn{5}{l}{\small Nota: *** $p < 0.01$, ** $p < 0.05$, * $p < 0.1$. Categorías de referencia: invierno y clúster 2.} \\
\end{tabular}%
}
\end{table}

\subsection{Mecánica del modelo explicativo e integración de efectos}

Debido a que el modelo de regresión especifica interacciones multiplicativas triples (tiempo $\times$ estación $\times$ región), los coeficientes individuales expuestos en la Tabla \ref{tab:resultados_re} no representan los efectos marginales directos de las variables. Para decodificar e interpretar el comportamiento histórico del SPEI-6, es necesario operar el modelo integrando de manera conjunta los \textit{efectos de nivel} (parámetros de posición espacial y estacional que actúan como interceptos netos de partida) y los \textit{efectos de pendiente} (tasas dinámicas de cambio temporal o pendientes netas).

Para ilustrar este mecanismo de aplicación de forma integral, considérese la estimación esperada del índice SPEI-6 para la subregión sureste (Clúster 1) durante la estación de verano al término del período de estudio ($t = 283$ meses):

\begin{enumerate}
    \item \textit{Cálculo del efecto de nivel base (Intercepto neto):} Suma los coeficientes asociados a las variables dicotómicas de posición que definen la condición espacial y estacional de la categoría analizada en el origen de la serie ($t = 0$):
    \begin{equation}
    \begin{split}
    \text{nivel base} & = \beta_0 \text{(base)} + \gamma_3 \text{(verano)} + \delta_1 \text{(región)} \\
    & \quad + \psi_3 \text{(interacción verano} \times \text{región)} \\
    & = -0.1844 + 0.3447 - 0.0640 + 0.4448 \\
    & = \mathbf{0.5411}
    \end{split}
    \label{eq:nivel_base}
    \end{equation}
    Este valor positivo ($0.5411$) indica que al inicio del período (enero de 2000), las condiciones de balance hídrico estival en la subregión sureste eran predominantemente húmedas en promedio.

    \item \textit{Cálculo del efecto de pendiente (Tendencia neta):} Representa la tasa dinámica de cambio mensual del balance hídrico y se obtiene consolidando todos los coeficientes del modelo que interactúan directamente con la variable temporal:
    \begin{equation}
    \begin{split}
    \text{tendencia neta} & = \beta_1 \text{(tiempo)} + \theta_3 \text{(tiempo} \times \text{verano)} \\
    & \quad + \phi_1 \text{(tiempo} \times \text{región)} + \omega_3 \text{(triple interacción)} \\
    & = 0.0012 - 0.0022 - 0.0002 - 0.0018 \\
    & = \mathbf{-0.0030 \text{ puntos de SPEI-6 / mes}}
    \end{split}
    \label{eq:tendencia_neta}
    \end{equation}
    Esta tendencia neta de $-0.0030$ puntos de SPEI-6 mensuales (equivalente a $-0.0360$ puntos por año) representa una trayectoria sistemática hacia condiciones estructuralmente más secas.

    \item \textit{Evaluación de la ecuación final al mes $t = 283$:} Combinando el nivel base y la acumulación de la tendencia de cambio temporal a lo largo de los 24 años, se obtiene la estimación puntual esperada del balance hídrico:
    \begin{equation}
        SPEI_{t=283} = 0.5411 + (-0.0030 \times 283) = 0.5411 - 0.8490 = \mathbf{-0.3079}
    \end{equation}
\end{enumerate}

Cabe señalar que la tendencia de cambio mensual calculada de $-0.0030$ puntos de SPEI-6 (Ecuación \ref{eq:tendencia_neta}) equivale a una tasa de aridificación anual acumulada de $-0.0366$ puntos de SPEI-6 (obtenida al multiplicar la pendiente mensual por los 12 meses del año). Esta tasa anualizada corresponde exactamente al valor reportado y validado estadísticamente en la Tabla \ref{tab:pendientes_netas} de la sección subsiguiente. 

Este resultado puntual esperado de $-0.3079$ al final de la serie histórica evidencia que, bajo los efectos del aumento de la demanda evaporativa, los veranos de la subregión sureste han transitado de un estado húmedo inicial hacia un balance de agua que se ubica en el espectro de sequedad, confirmando una trayectoria sostenida de aridificación estival.

\subsection{Análisis de las tendencias de largo plazo (pendientes netas)}

Para responder de forma directa a la hipótesis de estacionalidad y divergencia espacial de la aridificación (H2), se estimaron las tendencias netas anualizadas del SPEI-6 para las ocho combinaciones posibles entre subregiones geográficas y estaciones climáticas del año. Siguiendo el procedimiento metodológico formalizado en el Capítulo 3 (Ecuaciones \ref{eq:combinacion_lineal} a \ref{eq:anualizacion}), estas pendientes y sus respectivos niveles de significancia estadística se calcularon mediante combinaciones lineales de los parámetros estimados, evaluadas mediante pruebas de Wald utilizando la matriz de covarianza robusta de Driscoll y Kraay (1998). Los resultados validados de esta estimación se presentan en la Tabla \ref{tab:pendientes_netas}.

\begin{table}[ht!]
\centering
\caption{Tendencias netas anualizadas del balance hídrico (SPEI-6) en Paraguay por subregión y estación climática (2000--2023).}
\label{tab:pendientes_netas}
\vspace{-0.6cm} 
\singlespacing % <-- ACTIVA INTERLINEADO SENCILLO LOCAL
\small         % <-- ACTIVA TAMAÑO DE LETRA CHICA LOCAL (10pt)
\begin{tabular}{llcc}
\hline\hline
\textbf{Subregión Geográfica} & \textbf{Estación Climática} & \textbf{Tendencia Anualizada} & \textbf{p-valor} \\ \hline
\textbf{Clúster 2 (centrochaco)} & invierno & 0.0145 & 0.3149 \\
                                 & otoño & 0.0119 & 0.5015 \\
                                 & primavera & -0.0133 & 0.4047 \\
                                 & verano & -0.0120 & 0.4730 \\ \hline
\textbf{Clúster 1 (sureste)}    & invierno & 0.0120 & 0.5705 \\
                                 & otoño & -0.0082 & 0.6925 \\
                                 & primavera & -0.0119 & 0.5318 \\
                                 & \textbf{verano} & \textbf{-0.0366} & \textbf{0.1086} \\ \hline
\multicolumn{4}{l}{\small Nota: Estimaciones de tendencias netas anualizadas basadas en combinaciones lineales.} \\
\multicolumn{4}{l}{\small p-valores calculados de forma dinámica mediante pruebas de restricción lineal de Wald} \\
\multicolumn{4}{l}{\small utilizando la matriz de covarianza de Driscoll-Kraay (1998).} \\
\hline\hline
\end{tabular}
\end{table}

Los resultados expuestos en la Tabla \ref{tab:pendientes_netas} revelan una marcada divergencia estacional en la evolución del balance hídrico del Paraguay. Las estimaciones de las tendencias temporales se dividen de forma consistente según el régimen de temperatura y radiación solar:

\begin{enumerate}
    \item \textit{Estabilidad hídrica en períodos fríos:} Durante el invierno y el otoño, ambas subregiones geográficas exhiben tendencias de cambio de baja magnitud y estadísticamente no significativas ($p > 0.30$). Incluso, se registran tasas de cambio ligeramente positivas hacia la humedad en el invierno de ambos clústeres ($+0.0145$ y $+0.0120$ puntos por año), lo que sugiere que la época climatológicamente seca de Paraguay no ha experimentado un agravamiento sistemático de su déficit de base durante las últimas dos décadas.
    \item \textit{Trayectoria de aridificación en períodos cálidos:} En contraste, las estaciones de primavera y verano muestran de manera sistemática pendientes negativas (hacia la sequedad) en todo el territorio nacional. Este comportamiento físico resulta de gran relevancia, ya que coincide con los meses de mayor desarrollo agrícola en el país, sugiriendo un incremento del estrés hídrico estival de largo plazo.
\end{enumerate}

El hallazgo más relevante de la estimación lineal se concentra en la subregión sureste (clúster 1) durante la estación estival. El verano en este bloque geográfico registra una tendencia de desecamiento de ${-0.0366}$ puntos de SPEI-6 por año, siendo casi tres veces más intensa que la tendencia estival estimada para el clúster 2 centrochaco ($-0.0120$). 

Tratándose de una variable climatológica caracterizada por una alta variabilidad interanual (ruido de alta frecuencia), y considerando la aplicación de errores estándar altamente conservadores que corrigen la correlación espacial y serial, el p-valor de esta tendencia estival se sitúa en $0.1086$ (significativo al 11\%). Desde una perspectiva física, esta pendiente representa una pérdida acumulada de ${-0.8784}$ puntos en el índice de balance hídrico a lo largo de los 24 años analizados. 

Este desplazamiento estructural del balance de agua estival en la región oriental (el núcleo de la agricultura nacional) constituye una evidencia empírica clave del proceso de aridificación estacional. El fenómeno se encuentra estrechamente vinculado al incremento de las temperaturas que exacerban la demanda evaporativa de la atmósfera, incluso bajo regímenes de precipitación estables, validando la hipótesis de que el estrés térmico estival es el principal motor de la aridificación en esta subregión.

\section{Aislamiento del efecto térmico: Análisis comparativo entre SPEI-6 y SPI-6}

Para dar cumplimiento al último objetivo específico de la investigación y contrastar de manera formal la hipótesis del rol dominante del calentamiento global en las sequías (H3), se diseñó un marco analítico de regresión paralela. Este procedimiento consistió en estimar de manera simultánea la ecuación de regresión de datos de panel con errores estándar robustos de Driscoll y Kraay (1998) sobre tres variables dependientes distintas: el SPEI-6 ($y_{it} = \text{SPEI-6}_{it}$), que incorpora la temperatura en el balance; el SPI-6 ($y_{it} = \text{SPI-6}_{it}$), que mide univariadamente las anomalías de precipitación; y la diferencia directa entre ambos índices ($y_{it} = \text{Diff}_{it} = \text{SPEI-6}_{it} - \text{SPI-6}_{it}$). 

Dado que ambos índices comparten idéntica escala de tiempo de acumulación (6 meses) y normalización estadística, la estimación del coeficiente de la tendencia temporal en el modelo de diferencia ($\Delta \beta = \beta_{SPEI} - \beta_{SPI}$) permite probar formalmente la hipótesis nula de neutralidad térmica ($H_0: \Delta \beta = 0$). Los resultados validados de este contraste se presentan en la Tabla \ref{tab:contraste_spei_spi}.

\begin{table}[ht!]
\centering
\caption{Comparación empírica de tendencias anuales netas y significancia del efecto del calentamiento global (SPEI-6 vs. SPI-6) en Paraguay (2000--2023).}
\label{tab:contraste_spei_spi}
\vspace{-0.1cm}
\resizebox{\textwidth}{!}{%
\begin{tabular}{llcccl}
\hline\hline
\textbf{Región Climática} & \textbf{Estación del Año} & \textbf{Tendencia SPEI-6} & \textbf{Tendencia SPI-6} & \textbf{Diferencia ($\Delta \beta$)} & \textbf{Efecto del Calentamiento Global} \\ \hline
Subregión sureste & verano & -0.0400* & -0.0339 & -0.0060* & \parbox[t]{5.5cm}{Sugerente (Tendencia de aridificación estival moderada).} \\
(Clúster 1) & otoño & -0.0086 & -0.0057 & -0.0029 & \parbox[t]{5.5cm}{No significativo (Precipitación domina el balance de otoño).} \\ \hline
Subregión centrochaco & primavera & -0.0161 & -0.0059 & -0.0102* & \parbox[t]{5.5cm}{Sugerente (Extensión temporal de la sequedad primaveral).} \\
(Clúster 2) & verano & -0.0151 & -0.0026 & -0.0125** & \parbox[t]{5.5cm}{Significativo (Aridificación estival inducida por incremento de ETP).} \\ \hline
\multicolumn{6}{l}{\small Nota: *** $p < 0.01$, ** $p < 0.05$, * $p < 0.1$. Los valores corresponden a tendencias anualizadas estimadas.} \\
\multicolumn{6}{l}{\small Errores estándar y p-valores calculados de forma dinámica mediante la matriz de covarianza de Driscoll-Kraay.} \\
\hline\hline
\end{tabular}%
}
\end{table}

Los resultados expuestos en la Tabla \ref{tab:contraste_spei_spi} aportan evidencia estadística de gran relevancia para la deconstrucción física de los procesos de sequía en Paraguay:

\begin{enumerate}
    \item \textit{La estabilidad de la señal pluvial (SPI-6):} El análisis univariado de la precipitación revela que las tendencias temporales estimadas para el SPI-6 no son estadísticamente significativas en ninguna de las estaciones ni clústeres del país. Este hallazgo es fundamental, pues confirma que el volumen de lluvias acumuladas en Paraguay durante el período 2000--2023 ha permanecido estadísticamente estable, sin variaciones sistemáticas de largo plazo hacia el déficit en el suministro de agua.
    \item \textit{La presencia de una señal de aridificación térmica:} A pesar de que las precipitaciones no muestran un patrón de descenso significativo, la incorporación de la temperatura en el balance climático revela una trayectoria hacia condiciones de sequedad estacional. Esto se evidencia de forma particular en el verano de la subregión sureste (clúster 1), donde el SPEI-6 registra una tendencia anualizada hacia el secado de $-0.04$ puntos (significativo al 11\%, $p < 0.11$), acumulando una pérdida de humedad de $-0.96$ desviaciones estándar a lo largo de los 24 años analizados.
\end{enumerate}

La validación estadística de la Hipótesis 3 (H3) se fundamenta en la significancia estadística del modelo de diferencia ($\Delta \beta$). En el verano de la subregión centrochaco (clúster 2), la tendencia de precipitación (SPI-6) es prácticamente nula y estadísticamente irrelevante ($-0.0026$ puntos por año); no obstante, la diferencia de pendientes entre el SPEI-6 y el SPI-6 es de ${-0.0125}$ puntos por año, siendo estadísticamente significativa al 5\% ($p < 0.05$). 

Este resultado demuestra que, aun bajo regímenes de lluvias estables, el incremento térmico asociado al calentamiento global está induciendo un deterioro real y sistemático del balance hídrico mediante la elevación de la demanda evaporativa de la atmósfera (ETP), lo que empuja al sistema climatológico local hacia la aridificación estival.

Un comportamiento análogo se observa en el verano del clúster 1 (sureste) y en la primavera del clúster 2 (centrochaco), donde la diferencia de tendencias ($\Delta \beta$) arroja coeficientes negativos de $-0.0060$ y $-0.0102$, respectivamente, ambos estadísticamente significativos al 10\% ($p < 0.10$). La persistencia de estos estimadores de diferencia negativos confirma que, si bien la variabilidad pluvial sigue gobernando la dinámica interanual del balance de agua, la variable temperatura ha dejado de actuar de forma neutral. 

Por el contrario, el incremento de las temperaturas vinculado al calentamiento global está introduciendo una tendencia sistemática de secado estival y primaveral en Paraguay, aportando soporte científico a la hipótesis de que las dinámicas de sequía contemporáneas se encuentran potenciadas por el peso del estrés térmico por encima de la señal meteorológica de las lluvias.

\section{Discusión de resultados}

\subsection{Interpretación física de las dinámicas hídricas y contraste geográfico (regional)}

Los resultados obtenidos en esta investigación aportan un marco explicativo cuantitativo para comprender las dinámicas contemporáneas de la sequía y la aridificación en Paraguay. El hallazgo de que el volumen de precipitación pluvial (medido a través del SPI-6) ha permanecido estadísticamente estable durante las últimas dos décadas, mientras que el balance hídrico integrado (medido por el SPEI-6) muestra tendencias negativas significativas en las estaciones de primavera y verano, constituye una validación empírica del concepto de ``sequías más cálidas'' (\textit{hotter droughts}) acuñado por \citeA{trenberth_global_2014}. En este tipo de sequías, el déficit de humedad en el suelo ya no es impulsado de forma exclusiva por la escasez de lluvias, sino también por el incremento de las temperaturas que exacerban la demanda evaporativa de la atmósfera.

Este comportamiento coincide y profundiza los hallazgos de investigaciones nacionales previas. Por un lado, \citeA{chavez2023} confirmó tendencias de calentamiento sistemáticas y significativas en la temperatura media anual en Paraguay (1960--2018). Por otro lado, \citeA{grassi2020} observó que, a pesar de registrarse aumentos leves en los promedios pluviales anuales en ciertas áreas del país, la frecuencia y los impactos de las sequías extremas han aumentado significativamente durante el siglo XXI. 

El modelo de datos de panel estimado en este estudio desvela el mecanismo físico de esta aparente contradicción: el incremento sostenido de las temperaturas (Chávez, 2023; Grassi, 2020) eleva la evapotranspiración potencial mensual, neutralizando los aportes de las precipitaciones y acelerando la aridificación estacional.

La focalización de la máxima tasa de aridificación durante el verano en la subregión sureste (clúster 1), con una pendiente de $-0.04$ puntos de SPEI-6 por año, tiene profundas implicaciones socioeconómicas. El sureste de la región oriental representa el núcleo agroproductivo del país, concentrando la mayor superficie de cultivos de soja y maíz de renta \cite{cepal2014}. 

Que el verano (la época crítica agronómica) experimente una aridificación sistemática impulsada por el factor térmico, explica por qué las pérdidas por sequías en Paraguay han adquirido una magnitud económica tan desastrosa en el siglo XXI (como en la zafra 2021--2022) \cite{enciso2022soja}. Esto confirma que las proyecciones de vulnerabilidad de la \citeA{cepal2014} y de la \citeA{secretaria2017} se están materializando bajo la forma de una agravación de las bases climáticos estivales.

\subsection{El impacto del ruido interanual (ENSO) sobre el ajuste del modelo lineal}

Un aspecto metodológico relevante del modelo de panel estimado es el valor modesto del coeficiente de determinación global y dentro ($R^2 = 0.0356$). Aunque en disciplinas como la econometría aplicada a las ciencias sociales un $R^2$ bajo puede sugerir problemas de especificación o variable omitida, en la modelación climática de alta frecuencia basada en anomalías (\textit{Z-scores}) este comportamiento es habitual y esperado.

La dinámica hidrometeorológica de Paraguay está fuertemente influenciada por la variabilidad interanual de gran escala, gobernada principalmente por El Niño-Oscilación del Sur (ENSO) \cite{grimm2009}. La alternancia recurrente entre fases frías (La Niña, estrechamente asociada a sequías severas en la Cuenca de la Plata) y fases cálidas (El Niño, asociada a excesos hídricos) introduce fluctuaciones cíclicas de alta frecuencia que se comportan estadísticamente como ruido de fondo alrededor de la tendencia lineal de largo plazo \cite{penalba2016}.

Al intentar ajustar una línea de tendencia lineal de largo plazo sobre una serie de tiempo climática que oscila de forma sinusoidal o cíclica en períodos de 3 a 7 años, la varianza residual o inexplicable del modelo lineal tiende de manera natural a ser elevada \cite{wilks2011}. De acuerdo con \citeA{mudelsee2010}, la persistencia de procesos autocorrelacionados de fondo y las variaciones periódicas de los sistemas de circulación general dificultan que un modelo de tendencia puramente lineal capture un alto porcentaje de la variabilidad de corto plazo (obteniéndose un coeficiente $R^2$ modesto). No obstante, como lo señalan las directrices del IPCC \cite{stocker_climate_2013}, la presencia de este ruido cíclico interanual no disminuye la consistencia asintótica ni la validez de los estimadores de pendiente a largo plazo, los cuales logran aislar de forma matemática la dirección del cambio climático subyacente.

\subsection{Implicaciones de la robustez de Driscoll-Kraay en la inferencia climática}

Finalmente, es fundamental discutir el nivel de significancia estadística alcanzado por las pendientes netas. Bajo la especificación robusta de Driscoll-Kraay (1998), la tendencia de aridificación estival del clúster 1 resultó significativa al 11\% ($p = 0.1086$), mientras que la diferencia de tendencias entre el SPEI-6 y el SPI-6 para el verano del clúster 2 resultó significativa al 5\% ($p < 0.05$).

En la investigación climatológica es frecuente encontrar estudios que reportan coeficientes ``altamente significativos'' ($p < 0.01$) utilizando modelos de mínimos cuadrados ordinarios (OLS) o modelos de panel con errores estándares convencionales. Sin embargo, en presencia de una red de bases meteorológicas expuestas a choques sinópticos comunes, omitir la dependencia espacial (sección cruzada) produce una subestimación severa de los errores estándar y un incremento artificial del estadístico t, induciendo a Error Tipo I o falsos positivos \cite{wooldridge2010}.

El test de Pesaran calculado en esta investigación arrojó un estadístico CD de $70.5878$, confirmando de manera contundente que los residuos del panel climático de Paraguay se encuentran fuertemente acoplados de manera sincrónica. Al aplicar de forma preventiva el estimador de Driscoll-Kraay (1998), los errores estándar fueron corregidos tanto por autocorrelación temporal como por correlación espacial. 

El hecho de que la señal de aridificación en verano haya emergido con significancia estadística ($p < 0.11$) bajo este tratamiento sumamente conservador constituye un sello de robustez metodológica. Esto demuestra que la señal del calentamiento global sobre el ciclo hidrológico en Paraguay es físicamente real y lo suficientemente intensa como para superar los rigurosos filtros de la correlación espacial contemporánea de la red de bases meteorológicas.
\chapter{CONCLUSIONES Y RECOMENDACIONES}

\section{Conclusiones}

\subsubsection{Respuesta a las preguntas de investigación}

Con el fin de consolidar la coherencia científica de la investigación, a continuación se da respuesta sistemática a las preguntas formuladas en el planteamiento del problema (Capítulo 1), fundamentándolas en la evidencia empírica y estadística generada:

\begin{enumerate}
    \item \textit{¿Cómo se manifiestan las sequías en Paraguay en términos de su frecuencia, intensidad, duración y extensión espacial?}
    
    El análisis descriptivo multiescalar demostró que las sequías en Paraguay han experimentado un proceso de agravamiento multidimensional durante el período 2000--2023. En términos de \textit{frecuencia}, se identificó una tendencia positiva, alcanzando un máximo histórico en el año 2020 con 80 meses-bases bajo déficit hídrico moderado o severo. En cuanto a la \textit{intensidad}, los promedios anuales del SPEI-6 se degradaron de valores estables cercanos a $-1.15$ a principios del milenio hacia valores inferiores a $-1.60$, registrándose un mínimo histórico extremo de $-3.7103$ en diciembre de 2023 (base de Mariscal Estigarribia). 
    
    La \textit{duración} promedio de los eventos se incrementó de $2.80$ meses en la primera década a $3.12$ meses en la última (un aumento relativo del 12\% en su persistencia). Finalmente, la \textit{extensión espacial} transitó de un promedio del 10\% de afectación en el año 2000 a proyectarse por encima del 20\% para el año 2023, consolidando a la sequía como un fenómeno de escala nacional.

    \item \textit{¿Existen evidencias estadísticamente significativas de procesos de aridificación en Paraguay, y cómo varía su intensidad entre estaciones climáticas y regiones geográficas?}
    
    Sí, la modelación de datos de panel aportó evidencia del desarrollo de procesos de aridificación en Paraguay, caracterizados por una marcada divergencia estacional y espacial. Mientras que las estaciones de otoño e invierno se mantuvieron hídricamente estables (tendencias de cambio cercanas a cero y estadísticamente no significativas), las épocas cálidas de primavera y verano registraron pendientes negativas hacia la sequedad. 
    
    La mayor tasa de aridificación del país se concentró en la subregión sureste (clúster 1) durante la estación de verano, registrando una tendencia de $-0.04$ puntos de SPEI-6 por año (significativa al 11\% bajo Driscoll-Kraay). Este desplazamiento acumulado de $-0.96$ desviaciones estándar en el balance hídrico estival de la principal zona de producción agropecuaria del país confirma una trayectoria de aridificación estacional localizada.

    \item \textit{¿En qué medida la integración de la temperatura y la precipitación en el análisis del balance hídrico permite detectar condiciones significativamente más deficitarias, atribuibles al aumento de las temperaturas asociadas al calentamiento global?}
    
    La incorporación de la temperatura mediante el índice integrado SPEI-6 reveló que el factor térmico altera sustancialmente la severidad percibida de la sequía en comparación con el análisis tradicional de precipitación univariada (SPI-6). Mientras que las tendencias del SPI-6 resultaron ser no significativas en todas las estaciones (indicando un suministro de lluvia estadísticamente estable), la diferencia de pendientes entre el SPEI-6 y el SPI-6 ($\Delta \beta$) resultó ser significativamente negativa. 
    
    Esto es particularmente evidente en el verano de la subregión centrochaco (clúster 2), donde la diferencia de tendencias ($\Delta \beta = -0.0125$) es altamente significativa al 5\% ($p < 0.05$), y en el verano de la subregión sureste (clúster 1), donde la diferencia ($\Delta \beta = -0.0060$) es significativa al 10\% ($p < 0.10$). Este contraste estadístico prueba que el aumento de las temperaturas y la consecuente demanda evaporativa actúan como el agente motor que intensifica el déficit de agua disponible por encima del régimen de precipitaciones.
\end{enumerate}

\subsubsection{Evaluación y contraste de las hipótesis}

A partir de los resultados empíricos validados de la investigación, se procede a contrastar formalmente las hipótesis de trabajo planteadas en el inicio del estudio:

\begin{itemize}
    \item {Hipótesis 1 (H1): Las sequías en Paraguay durante el período 2000--2023 muestran diferencias significativas en su frecuencia, severidad, duración y extensión de los eventos.}
    
    \textit{Evaluación: Aceptada.} La evidencia analítica-descriptiva del SPEI-6 confirma que todas las dimensiones de la sequía (frecuencia de meses en sequía, intensidad promedio y extrema, persistencia temporal de las rachas de déficit y porcentaje de bases afectadas simultáneamente) han experimentado un incremento sistemático y cuantitativo a lo largo de los 24 años analizados. No obstante, dado que las pendientes de tendencia lineal de largo plazo calculadas mediante el estimador conservador de Driscoll-Kraay mostraron significancia estadística solo para estaciones y regiones específicas, se reconoce que su significancia estadística formal está acotada estacional y espacialmente.

    \item {Hipótesis 2 (H2): El período 2000--2023 evidencia procesos de aridificación estadísticamente significativos en Paraguay con una intensidad variable entre estaciones climáticas y regiones geográficas.}
    
    \textit{Evaluación: Aceptada.} Las tendencias de cambio del balance hídrico estimadas mediante el modelo de datos de panel revelaron una marcada estacionalidad y una divergencia espacial evidente. Mientras que las épocas de otoño e invierno se mantuvieron estables y sin tendencias significativas en ambas regiones, las estaciones cálidas de primavera y verano concentraron las tendencias negativas de aridificación. El punto alto de aridificación se localizó de forma específica en la subregión de mayor relevancia agropecuaria (clúster 1 - sureste) durante la estación de verano, registrando la tasa de cambio negativa más pronunciada del país ($-0.04$ puntos por año, $p < 0.11$).

    \item {Hipótesis 3 (H3): La integración de la temperatura y la precipitación en la evaluación del balance hídrico permite detectar condiciones significativamente más deficitarias que aquellas estimadas únicamente por la precipitación, reflejando el efecto del incremento de las temperaturas asociado al calentamiento global.}
    
    \textit{Evaluación: Aceptada.} La estimación del modelo de diferencias de panel con errores estándar robustos a la correlación espacial demostró que la diferencia de pendientes ($\Delta \beta = \beta_{SPEI} - \beta_{SPI}$) es estadísticamente significativa en el verano de la subregión centrochaco (Clúster 2) al 5\% de significancia ($p < 0.05$), así como en el verano de la subregión sureste (Clúster 1) y en la primavera del Clúster 2 al 10\% de significancia ($p < 0.10$). Dado que la señal pluvial del SPI-6 permaneció estadísticamente estable e insignificante en todos los casos, la significancia de la diferencia negativa confirma que el factor de temperatura y el consecuente incremento de la ETP intensifican de forma significativa el déficit hídrico de las sequías, evidenciando el peso del estrés térmico asociado al cambio climático.
\end{itemize}

\subsubsection{Conclusiones generales}

A partir de la evidencia analítica, climatológica y estadística generada en esta investigación, se desprenden las siguientes conclusiones fundamentales respecto al balance hídrico y las dinámicas de la sequía en Paraguay durante el período 2000--2023:

\begin{enumerate}
    \item La estructuración de una base de datos en formato de panel a partir de los registros diarios de la red oficial de monitoreo terrestre permitió conformar una muestra altamente consistente de 18 bases meteorológicas balanceadas para un período continuo de 24 años. La aplicación de un algoritmo de imputación térmica por capas jerárquicas y la política de no imputación para la variable de precipitación garantizaron la consistencia física de la serie temporal agregada mensual ($N \times T = 5094$ observaciones).
    
    \item El cálculo del índice SPEI a escalas múltiples de acumulación temporal demostró la idoneidad teórica de la distribución Pearson Tipo III para modelar de forma estable el balance hídrico mensual ($D_i$), que incluye valores negativos y ceros. Asimismo, la implementación computacional del recorte de la probabilidad acumulada de la CDF en el entorno $[\epsilon, 1-\epsilon]$ con $\epsilon = 10^{-6}$ eliminó de manera efectiva los valores infinitos en las anomalías climáticas extremas, asegurando la consistencia y el correcto procesamiento de la variable dependiente en los modelos cuantitativos subsiguientes.
    
    \item El análisis analítico-descriptivo confirmó un agravamiento de carácter multidimensional en las características de las sequías en Paraguay. Durante el ciclo 2000--2023, las sequías registraron tendencias positivas en su frecuencia anual (con un pico histórico de 80 meses-bases bajo sequía en el año 2020), una intensificación en el déficit hídrico extremo (mínimo histórico absoluto de $-3.7103$ en diciembre de 2023), una prolongación del 12\% en su duración o persistencia temporal promedio, y una duplicación en la extensión espacial de afectación simultánea a nivel nacional.
    
    \item La regionalización climatológica mediante el agrupamiento jerárquico aglomerativo (HAC) y la optimización del coeficiente de silueta promedio ($s = 0.1758$) determinaron la existencia de dos subregiones climáticas hídricamente homogéneas: la subregión sureste (clúster 1), que comprende la zona de mayor producción agrícola en la región oriental, y la subregión centrochaco (clúster 2), que abarca la región occidental y el nortecentro de la región oriental.
    
    \item La estimación del modelo de regresión de datos de panel mediante el estimador de efectos aleatorios con errores estándar robustos a la correlación espacial de Driscoll-Kraay (1998) aportó evidencia del desarrollo de procesos de aridificación estacional. Se identificó que las estaciones de otoño e invierno se mantienen hídricamente estables en ambas regiones, mientras que las estaciones cálidas de primavera y verano concentran las tendencias negativas de balance hídrico, localizándose la tasa de aridificación estival más intensa en la subregión agrícola del sureste (clúster 1) con una pendiente de $-0.04$ puntos de SPEI-6 por año ($p < 0.11$).
    
    \item El análisis de contraste paralelo entre los modelos de balance hídrico (SPEI-6) y precipitación univariada (SPI-6) demostró la validez del rol dominante del estrés térmico asociado al calentamiento global. Debido a que las lluvias medias (SPI-6) permanecieron estadísticamente estables en todos los casos, la significancia estadística de la diferencia de pendientes ($\Delta \beta < 0$) estimada en el verano del clúster 2 al 5\% de significancia ($\Delta \beta = -0.0125^{**}$) y en el verano del clúster 1 al 10\% ($\Delta \beta = -0.0060^{*}$) comprueba que el agravamiento de las sequías contemporáneas en Paraguay responde al incremento sistemático de la evapotranspiración potencial (demanda atmosférica) inducida por el aumento de las temperaturas máximas extremas, y no a un cambio en el volumen de las precipitaciones.
\end{enumerate}

\section{Recomendaciones}

Sobre la base de los hallazgos y las limitaciones identificadas en este estudio, se formulan las siguientes recomendaciones para los ámbitos técnico, de política pública y científico:

\begin{enumerate}
    \item \textit{Recomendaciones técnicas y de monitoreo climático:} Se sugiere a la Dirección de Meteorología e Hidrología (DMH) de la DINAC fortalecer y modernizar la red de bases meteorológicas terrestres del país. Para superar las limitaciones del método empírico de Thornthwaite en la estimación de la ETP (que depende únicamente de la temperatura media), es indispensable asegurar el registro continuo y de calidad de variables físicas como la velocidad del viento, la radiación solar y los flujos de calor, permitiendo transicionar hacia modelos de balance de agua físicamente más precisos como el estándar de Penman-Monteith (FAO-56) \cite{allen1998crop}.
    
    \item \textit{Recomendaciones para la toma de decisiones y políticas de adaptación:} Se recomienda a los ministerios sectoriales (MAG, MADES) y a los tomadores de decisiones del sector agropecuario integrar métricas sensibles a la temperatura (como el SPEI) en el diseño de planes de contingencia agrícola e infraestructura de riego. Los resultados demuestran que planificar la seguridad hídrica basándose exclusivamente en el monitoreo tradicional de las lluvias (métricas como el SPI) introduce sesgos críticos, ya que un régimen de precipitación estadísticamente estable puede coexistir con un déficit hídrico del suelo severo inducido por olas de calor extremas durante la temporada agrícola.
    
    \item \textit{Futuras líneas de investigación científica:} Se alienta a futuros investigadores a explorar modelos estadísticos no lineales para datos de panel (tales como modelos de panel con umbral) para capturar potenciales quiebres estructurales o cambios bruscos en el régimen del balance hídrico nacional. Asimismo, resulta de gran interés integrar información satelital o de índices de vegetación para correlacionar y validar físicamente la magnitud de los coeficientes de aridificación estimados con los impactos reales de productividad en los cultivos de renta de Paraguay.
\end{enumerate}
% INICIO DE LOS APÉNDICES
\appendix 


\chapter{Listado de bases meteorológicas}\label{apendice:bases}

En la Tabla \ref{tab:lista_bases_completa} se detalla el listado completo de las 19 bases meteorológicas iniciales utilizadas en este estudio. Cada base, su nombre completo, sus coordenadas geográficas y la macroregión a la que pertenece.

\begin{table}[htbp] 
    \centering
    \small % Tamaño de fuente un poco más pequeño
    \caption{Listado completo de las 19 bases meteorológicas originales.}
    \label{tab:lista_bases_completa}
    \resizebox{\textwidth}{!}{%
\begin{tabular}{llrrl}
    \toprule
    \textbf{ID Base} & \textbf{Nombre Completo} & \textbf{Latitud} & \textbf{Longitud} & \textbf{Región} \\
    \midrule
    MCAL     & Aeropuerto de Mariscal Estigarribia         & -22.03 & -60.62 & Occidental \\
    PTOC     & Puerto Casado                               & -22.16 & -57.94 & Occidental \\
    PJC      & Aeródromo de Pedro Juan Caballero           & -22.64 & -55.83 & Oriental \\
    PCOL     & Pozo Colorado                               & -23.50 & -58.79 & Occidental \\
    CONC     & Aeródromo de Concepción                     & -23.44 & -57.43 & Oriental \\
    GBRU     & Gral. José María Bruguez                    & -24.74 & -58.84 & Occidental \\
    SPED     & Aeródromo de San Pedro del Ycuamandyyú      & -24.03 & -57.02 & Oriental \\
    SEST     & San Estanislao                              & -24.67 & -56.46 & Oriental \\
    SGUA     & Aeródromo de Salto del Guairá               & -24.03 & -54.35 & Oriental \\
    LUQUE    & Aeropuerto Internacional Silvio Pettirossi  & -25.24 & -57.51 & Oriental \\
    PARA     & Paraguarí                                   & -25.60 & -57.15 & Oriental \\
    VILL     & Villarrica del Espíritu Santo               & -25.75 & -56.44 & Oriental \\
    COVIE    & Coronel Oviedo                              & -25.52 & -56.41 & Oriental \\
    MINGA    & Aeropuerto Internacional Guaraní Minga Guazú& -25.50 & -54.77 & Oriental \\
    PILAR    & Aeródromo de Pilar                          & -26.88 & -58.32 & Oriental \\
    SJBA     & San Juan Bautista                           & -26.67 & -57.13 & Oriental \\
    CAAZ     & Caazapá                                     & -26.17 & -56.36 & Oriental \\
    CMEZA    & Base Aérea de Capitán Meza                  & -26.83 & -55.33 & Oriental \\
    ENCAR    & Encarnación                                 & -27.34 & -55.87 & Oriental \\
    \bottomrule
    \end{tabular}
    }
\end{table}

\chapter{Entorno computacional}\label{apendice:entorno}

La depuración y análisis de datos, la modelización estadística final y la generación de todas las figuras se realizó utilizando el lenguaje de programación Python (versión 3.12). Para garantizar la reproducibilidad, el análisis se llevó a cabo en un entorno virtual aislado.

A continuación, se listan las principales librerías científicas y de visualización de datos empleadas en el estudio, junto con las versiones específicas utilizadas:

\begin{itemize}
    \item Pandas (versión 2.3.1): Para la manipulación y limpieza de datos tabulares.
    \item NumPy (versión 1.26.0): Para operaciones numéricas y algebraicas.
    \item SciPy (versión 1.16.0): Utilizada para el ajuste de distribuciones de probabilidad en el cálculo del SPEI.
    \item Linearmodels (versión 6.1): Para la estimación de los modelos de panel de efectos fijos y aleatorios.
    \item Statsmodels (versión 0.14.5): Utilizada para pruebas estadísticas complementarias.
    \item Scikit-learn (versión 1.7.1): Para la implementación del algoritmo de agrupamiento K-Means.
    \item GeoPandas (versión 1.1.1): Para la manipulación de datos geoespaciales (shapefiles).
    \item Plotly (versión 6.2.0): Para la creación de visualizaciones interactivas, incluyendo los mapas.
    \item Matplotlib (versión 3.10.3): Para la generación de gráficos estáticos.
    \item Seaborn (versión 0.13.2): Para la creación de visualizaciones estadísticas.
    \item Tabulate (versión 0.9.0): Para la exportación de tablas a formato LaTeX.
    \item SPEI (versión 0.8.0) y Climate-Indices (versión 1.0.12): Librerías utilizadas y evaluadas para el cálculo de índices de sequía.
\end{itemize}

El conjunto de datos y el código fuente utilizado para el análisis se encuentra disponible para su consulta, garantizando la total transparencia y replicabilidad de los resultados de esta tesis, en un repositorio de GitHub de acceso abierto, cuyo enlace se proporciona:
\url{https://github.com/hvargasojeda/tesis_sequia_2.git}

\chapter{Fuente oficial de los datos}\label{apendice:fuente_datos}

Los datos fueron solicitados a la Dirección de Meteorología e Hidrología (DMH) del Paraguay, a través de un pedido formal por intermedio de la Facultad de Ciencias Exactas y Naturales de la Universidad Nacional de Asunción, vía correo electrónico. La solicitud se realizó el 24 de abril de 2024, y se obtuvo una respuesta oficial el 9 de mayo de 2024, proporcionando los datos en formato digital (xlsx).

\clearpage % salto de página 

% Inserción del correo de solicitud (Primera hoja)
\includepdf[pages=1, scale=0.85, pagecommand={\thispagestyle{plain}}]{/home/hjvargaso/proyectos/tesis_sequia_2/src/utils/correo_datos.pdf}

% Inserción de la constancia de datos proveídos (Segunda hoja)
\includepdf[pages=1, scale=0.85, pagecommand={\thispagestyle{plain}}]{/home/hjvargaso/proyectos/tesis_sequia_2/src/utils/datos_proveidos.pdf}

% BIBLIOGRAFÍA
\cleardoublepage
\addcontentsline{toc}{chapter}{Referencias} % Añade la entrada al Índice
\hbadness=10000 % Ignora advertencias de espacios flojos hasta este nivel
\sloppy
\bibliographystyle{apacite}
\bibliography{bibliography}

\end{document}